%=====================================================================
%  THE AI AND DATA SCIENCE COURSE CLUSTER ACROSS THREE HEC CURRICULA
%
%  Compares HEC 2017, 2023 and 2025 for one group of courses only:
%  AI, machine learning, data mining, data science, deep learning, NLP
%  and their neighbours. Ends with a pathway and semester-wise schemes
%  for BSIT and for BSCS with an IT specialization.
%
%  Build:  latexmk -pdf HEC_AI_DS_Curriculum_Map.tex
%     or:  .\build.ps1
%=====================================================================
\documentclass[11pt,a4paper,oneside]{report}
\usepackage{hecmap}

\begin{document}

%---------------------------------------------------------------------
% Title page
%---------------------------------------------------------------------
\thispagestyle{empty}
\begin{center}
\vspace*{1.0cm}

{\color{secondary}\rule{0.85\textwidth}{2pt}}\\[6pt]
{\large\mdseries\textsc{\color{secondary}Higher Education Commission of Pakistan}}\\[16pt]

{\Huge\bfseries\color{primary}THE AI AND DATA SCIENCE}\\[6pt]
{\Huge\bfseries\color{primary}COURSE CLUSTER}\\[10pt]
{\LARGE\bfseries\color{primary}Across Three Curriculum Editions}\\[16pt]

{\color{secondary}\rule{0.85\textwidth}{2pt}}\\[16pt]

{\large\mdseries\textsc{2017 \;$\bullet$\; 2023 \;$\bullet$\; 2025}}\\[6pt]
{\small\color{gray!70!black}A version-by-version comparison, with a pathway and a
semester-wise scheme for BSIT and BSCS (IT specialization)}\\[24pt]

\begin{tikzpicture}[
  ed/.style={rounded corners=4pt, draw=#1, line width=1.1pt, fill=#1!6,
             text=#1, align=center, font=\small\bfseries,
             minimum width=4.4cm, minimum height=2.0cm}]
  \node[ed=ed2017] at (-4.7,0)
    {\faIcon{archive}\\[3pt]2017\\[2pt]\footnotesize\mdseries BS \& MS\\\footnotesize\mdseries CS, SE \& IT};
  \node[ed=ed2023] at (0,0)
    {\faIcon{sitemap}\\[3pt]2023\\[2pt]\footnotesize\mdseries BS Computing\\\footnotesize\mdseries Disciplines};
  \node[ed=ed2025] at (4.7,0)
    {\faIcon{rocket}\\[3pt]2025\\[2pt]\footnotesize\mdseries BS Computer\\\footnotesize\mdseries Science};
  \draw[-{Stealth[length=3mm]}, line width=1pt, gray!60] (-2.4,0) -- (-2.15,0);
  \draw[-{Stealth[length=3mm]}, line width=1pt, gray!60] (2.3,0) -- (2.55,0);
\end{tikzpicture}\\[14pt]
{\small\itshape\color{gray!70!black}Every claim in this document is taken from the
three source booklets and cited to them.}

\vspace{1.1cm}

\begin{tcolorbox}[enhanced, width=0.88\textwidth, colback=lightbg,
  colframe=primary, boxrule=0pt, leftrule=4pt, arc=2pt, top=6pt, bottom=6pt]
\small
\textbf{\color{primary}Scope.} This is not a review of the three curricula. It
covers one group of courses only --- artificial intelligence, machine learning,
data mining, data science, deep learning, natural language processing and the
courses that sit alongside them --- together with the three infrastructure
courses that the group runs on and competes with for credit: \emph{Computer
Networks}, \emph{System \& Network Administration} and \emph{Cloud Computing}.
Everything else in the three booklets (software engineering, cyber security,
graphics, the general education block) appears only where it is a prerequisite,
a competitor for elective credit, or a source of overlap.
\end{tcolorbox}

\vspace{0.9cm}
{\small\color{gray!60!black}Prepared for curriculum planning
$\bullet$ \today}
\end{center}

\newpage
\pagenumbering{roman}
\tableofcontents
\newpage
\pagenumbering{arabic}
\setcounter{page}{1}

%=====================================================================
\section{The Three Sources, and What Each One Is}
%=====================================================================

The three documents are not three drafts of the same thing. They differ in what
they cover, and that difference explains most of what follows.

\begin{longtable}{@{}L{1.9cm}L{4.5cm}L{4.4cm}L{4.5cm}@{}}
\toprule
\rowcolor{primary!15}
\textbf{} & \textbf{\color{ed2017}2017} & \textbf{\color{ed2023}2023} & \textbf{\color{ed2025}2025} \\
\midrule
\endhead
\textbf{Full title} &
BS \& MS (CS, SE \& IT) Curriculum &
BS Curriculum, Computing Disciplines &
HEC Curriculum of Computer Science \\
\rowcolor{lightbg}
\textbf{Covers} &
BSCS, BSSE, BSIT, and the MS programmes &
Eleven BS degrees, including BSAI, BSDS and BSIT, plus an associate degree &
BS Computer Science only, plus an associate degree in computing \\
\textbf{Length} & 171 pages & 100 pages & 86 pages \\
\rowcolor{lightbg}
\textbf{Degree size} & 130 credit hours & 130 credit hours & 130 credit hours \\
\textbf{Is there a BSIT?} &
Yes, a full separate programme &
Yes, a full separate programme &
No. IT survives only as one of fourteen specializations inside BSCS \\
\rowcolor{lightbg}
\textbf{How courses are specified} &
Credit hours, prerequisites, CLOs with Bloom levels, a full topic outline, reference books &
The same, plus a course introduction paragraph &
CLOs with Bloom levels only. No topic outline, no prerequisites, no reference books \\
\textbf{How the cluster is organised} &
A flat list of 75 bachelor courses; universities pick electives from it &
Each degree gets a named domain core and a domain elective list &
Fourteen named specialization clusters of 15--21 courses each \\
\bottomrule
\end{longtable}

\begin{finding}[The single most important structural fact]
In \yr{2017} the AI and data science cluster is essentially a \emph{postgraduate}
subject. Machine learning, deep learning, data mining and data science do not
exist as undergraduate courses at all --- they appear only inside MS Data
Science. By \yr{2023} the cluster has moved down to the bachelor level so
completely that it supports two whole degrees, BSAI and BSDS. By \yr{2025} it
has become a named 24-credit specialization that any BS Computer Science student
may elect. Three editions, three completely different answers to ``when should a
student meet machine learning?''
\end{finding}

\subsection{Where the elective credit actually is}

For a BSIT or BSCS student the cluster is almost never compulsory. It is bought
with elective credit, so the size and placement of the elective block decides how
much of it a student can take.

\begin{longtable}{@{}L{2.5cm}L{4.3cm}L{4.3cm}L{4.2cm}@{}}
\toprule
\rowcolor{primary!15}
\textbf{} & \textbf{\color{ed2017}2017 (BSIT)} & \textbf{\color{ed2023}2023 (BSIT)} & \textbf{\color{ed2025}2025 (BSCS)} \\
\midrule
\endhead
\textbf{Elective budget} &
15 cr, 5 courses (``IT Elective 1--5'') &
21 cr, 7 courses (Domain Electives) &
24 cr, 8 courses (Elective I--VIII) \\
\rowcolor{lightbg}
\textbf{Where they fall} &
Semesters 5--8 &
Semester 5 ($\times$2), Semester 6 ($\times$4), Semester 7 ($\times$1) &
Semester VI ($\times$4), Semester VII ($\times$4) \\
\textbf{Cluster courses available to spend it on} &
Five: AI, Big Data Analytics, Computer Vision, Digital Image Processing, NLP &
All of them --- the master list of BS courses carries every course used by BSAI and BSDS &
Fifteen in the Data Science cluster, fifteen in the AI cluster \\
\rowcolor{lightbg}
\textbf{Constraint} &
No ML, DL, DM or DS course exists to elect &
The BSIT elective examples name none of the cluster; a department must reach into the master list deliberately &
Choosing the named \emph{Information Technology} specialization spends all eight electives on courses with no AI or data content \\
\bottomrule
\end{longtable}

%=====================================================================
\section{Where Each Course Lives in Each Edition}
%=====================================================================

This is the whole comparison on one page. Read a row to see what happened to a
course over eight years; read a column to see what a student in that edition
could actually take.

\begin{center}
\small
\stcore\ Compulsory\quad
\stelec\ Available as an elective\quad
\stpg\ Postgraduate only\quad
\stabsent\ Does not exist
\end{center}

\begin{longtable}{@{}L{4.5cm}C{1.45cm}C{1.45cm}C{1.45cm}L{6.0cm}@{}}
\toprule
\rowcolor{primary!15}
\textbf{Course} & \textbf{2017} & \textbf{2023} & \textbf{2025} & \textbf{Note} \\
\midrule
\endhead
Artificial Intelligence & \stelec & \stcore & \stcore &
Compulsory for \emph{every} computing degree from 2023. In 2017 it was a BSCS
domain core course and entirely absent from BSIT \\
\rowcolor{lightbg}
Machine Learning & \stpg & \stelec & \stelec &
The pivotal change. MS-only in 2017; a BSAI domain core and BSDS elective in
2023; in both the DS and AI clusters in 2025 \\
Deep Learning & \stpg & \stelec & \stelec &
2023 splits it in two --- \emph{ANN \& Deep Learning} and \emph{Deep Learning};
2025 merges them back into one course \\
\rowcolor{lightbg}
Data Mining & \stabsent & \stelec & \stelec &
Absent from 2017 at every level. A BSDS domain core in 2023 \\
Introduction to Data Science & \stpg & \stelec & \stelec &
2017 has \emph{Tools and Techniques in Data Science} at MS level only \\
\rowcolor{lightbg}
Natural Language Processing & \stelec & \stelec & \stelec &
Present in all three, but the content is unrecognisable between 2017 and 2025 \\
Big Data Analytics & \stelec & \stelec & \stelec &
The most stable course in the cluster \\
\rowcolor{lightbg}
Computer Vision & \stelec & \stelec & \stelec &
A BSAI domain core in 2023 \\
Data Visualization & \stpg & \stelec & \stelec &
A BSDS domain core in 2023 \\
\rowcolor{lightbg}
Data Warehousing \& BI & \stabsent & \stelec & \stelec &
Renamed to plain \emph{Business Intelligence} in 2025, with warehousing moved
into \emph{Data Engineering} \\
Reinforcement Learning & \stabsent & \stelec & \stelec &
2017 covers it as a topic inside AI, never as a course \\
\rowcolor{lightbg}
Knowledge Representation \& Reasoning & \stabsent & \stelec & \stelec &
A BSAI domain core in 2023 \\
Programming for AI & \stabsent & \stelec & \stelec &
A BSAI domain core in 2023 \\
\rowcolor{lightbg}
Speech Processing & \stabsent & \stelec & \stelec & \\
Text Mining & \stabsent & \stelec & \stabsent &
Dropped in 2025; its ground is split between NLP and Information Retrieval \\
\rowcolor{lightbg}
Advanced / Computational Statistics & \stabsent & \stelec & \stelec &
Renamed \emph{Computational Statistics} in 2025 \\
Generative AI & \stabsent & \stabsent & \stelec &
New in 2025, and in both the DS and AI clusters \\
\rowcolor{lightbg}
Agentic AI & \stabsent & \stabsent & \stelec & New in 2025 \\
MLOps & \stabsent & \stabsent & \stelec & New in 2025 \\
\rowcolor{lightbg}
Data Engineering & \stabsent & \stabsent & \stelec & New in 2025 \\
Information Retrieval & \stabsent & \stabsent & \stelec &
New as a course; 2017 taught it inside NLP \\
\rowcolor{lightbg}
Data Ethics \& Security & \stabsent & \stabsent & \stelec & New in 2025 \\
Optimization Techniques & \stabsent & \stabsent & \stelec & New in 2025 \\
\rowcolor{lightbg}
Stochastic Processes & \stelec & \stabsent & \stelec &
In the 2017 bachelor list, gone in 2023, back in the 2025 AI cluster \\
Digital Image Processing & \stelec & \stelec & \stabsent &
Absorbed into Computer Vision by 2025 \\
\midrule
\multicolumn{5}{@{}l}{\textbf{\color{primary}The infrastructure the cluster runs on}}\\
\midrule
Computer Networks & \stcore & \stcore & \stcore &
Compulsory in every degree in every edition. The only course in this document
that is \emph{never} an elective \\
\rowcolor{lightbg}
System \& Network Administration & \stcore & \stcore & \stelec &
BSIT-compulsory in 2017 and 2023; demoted to one elective among fifteen in the
2025 IT specialization \\
Cloud Computing & \stabsent & \stelec & \stcore &
The mirror image: named but never specified in 2017, an unspecified BSIT elective
in 2023, a compulsory major course for every BSCS student in 2025 \\
\rowcolor{lightbg}
Virtual Systems \& Services & \stelec & \stelec & \stabsent &
2017's nearest thing to a cloud course --- a hypervisor course. Gone by 2025,
its ground taken by Cloud Computing \\
IT Infrastructure & \stcore & \stcore & \stelec &
The same trajectory as System \& Network Administration \\
\rowcolor{lightbg}
Network Security & \stabsent & \stelec & \stelec &
In three 2025 clusters at once: Cyber Security, Information Technology, and
Network Infrastructure \& Cloud Computing \\
\bottomrule
\end{longtable}

\begin{finding}[Why three networking courses belong in a document about machine learning]
Three reasons, and none of them is padding. \textbf{They are the platform.} Big
Data Analytics is a Spark course, MLOps is a deployment course and Data
Engineering is a pipeline course --- all three assume the machine underneath.
\textbf{They are the competition.} A BSIT student has seven elective slots; every
one spent on Cloud Computing is one not spent on Deep Learning, and Part~8 prices
that trade-off explicitly. \textbf{They are where HEC 2025 quietly put the
bridge:} the \emph{Network Infrastructure \& Cloud Computing} specialization
contains both \emph{Machine Learning} and \emph{MLOps}, which makes it the only
2025 cluster outside Data Science and AI to reach into this material at all.
\end{finding}

\begin{warnbox}[A BSIT student in 2017 could not study machine learning]
This is worth stating plainly because it explains why so many Pakistani IT
departments still have no ML teaching capacity. Under the 2017 curriculum the
BSIT computing core was ten courses --- programming, data structures, discrete
structures, operating systems, databases, software engineering, networks,
information security and the project --- and \emph{artificial intelligence was
not among them}. It was a BSCS domain core course only. The 75-course bachelor
list offered AI, Big Data Analytics, Computer Vision, Digital Image Processing
and NLP; it contained no course called Machine Learning, Deep Learning, Data
Mining or Data Science. A department that wanted to teach machine learning had
to invent the course itself.
\end{warnbox}

%=====================================================================
\section{Course by Course, Version by Version}
%=====================================================================

Each card below takes one course and shows what each edition says about it.
Credit-hour notation follows the source: 2017 writes \texttt{3-1} or
\texttt{3(2,1)}, 2023 writes \texttt{3 (2-3)} meaning two lecture hours and three
laboratory hours per week, and 2025 writes \texttt{3+0} or \texttt{2+1}.

These cards are the \emph{analysis}. The full published text --- every topic
outline, prerequisite, learning outcome and reading list, verbatim, for
thirty-three courses across all three editions --- is reproduced in Part~10,
which is the part to open when writing an actual course file.

\subsection{The AI and data science courses}

\begin{coursecard}{Artificial Intelligence}{the hinge the whole cluster turns on}
\edrow{2017}{BSCS domain core, 3+1. Absent from BSIT.}{Prerequisite
\emph{Data Structures and Algorithms}. Outline: basic components and branches of
AI; reasoning and knowledge representation, propositional and first-order logic;
uninformed, informed and local search; constraint satisfaction; adversarial
search with minimax and alpha--beta; learning (supervised, unsupervised,
reinforcement); uncertainty and fuzzy logic; recent trends. Russell \& Norvig.}
\edrow{2023}{Computing core for every degree, Semester 3, 3 (2-3).}{Prerequisite
\emph{Object Oriented Programming} --- notably weaker than 2017's data-structures
requirement. Outline drops the formal logic emphasis and adds case studies
(General Problem Solver, Eliza, Student, Macsyma), ANN and NLP as topics, and
mandates Python.}
\edrow{2025}{Major core, Semester V, 3 credit hours.}{No prerequisite published
and no topic outline --- only CLOs. The course now sits two years later in the
degree than it did in 2023.}
\verdict{It became compulsory for everyone, then slid from Semester 3 to
Semester V. The 2023 move is the single most consequential change in these three
documents for IT students; the 2025 move squeezes everything that depends on it
into the final two years.}
\end{coursecard}

\begin{coursecard}{Machine Learning}{the course that did not exist for undergraduates in 2017}
\edrow{2017}{MS Data Science programme core only. 3 credit hours.}{One of three
MS(DS) core courses, alongside \emph{Statistical and Mathematical Methods for
Data Science} and \emph{Tools and Techniques in Data Science}. There is no
undergraduate equivalent anywhere in the booklet.}
\edrow{2023}{BSAI domain core (Semester 4); BSDS domain elective. 3 (2-3).}{Prerequisite
\emph{Artificial Intelligence}. Five CLOs from C1 to C6. Outline is the classical
Mitchell syllabus: concept learning and version spaces; decision trees, naive
Bayes, ANNs, SVMs; overfitting, pruning, accuracy measurement; linear and
logistic regression; hierarchical and $k$-means clustering, SOM, $k$-NN;
semi-supervised EM; reinforcement learning with HMMs, Monte Carlo and MDPs;
ensembles, bagging and boosting.}
\edrow{2025}{In three clusters at once: Data Science, AI, and Software Engineering.}{3+0
in the DS and AI clusters, 2+1 in the Software Engineering cluster --- the same
course with different laboratory provision depending on which list you read it
from. CLOs add an explicit ethics outcome (A3: Valuing).}
\verdict{From postgraduate to ubiquitous in eight years. The 2023 outline is the
most detailed specification of the three and remains the best reference for what
the course should contain; 2025 gives it an ethics outcome but no syllabus.}
\end{coursecard}

\begin{coursecard}{Deep Learning}{one course, then two, then one again}
\edrow{2017}{MS Data Science specialization core. 3 credit hours.}{One of four
courses from which MS(DS) students choose two. \emph{Deep Reinforcement Learning}
also appears in the MS elective list. Nothing at bachelor level.}
\edrow{2023}{Two separate courses.}{\emph{Artificial Neural Networks \& Deep
Learning}, a BSAI domain core, prerequisite \emph{Programming for AI} --- covers
perceptron and Adaline, Hebbian learning, Kohonen maps, associative and
bidirectional associative memory, Boltzmann machines, backpropagation, CNNs,
autoencoders, deep belief networks. Then \emph{Deep Learning}, an AI elective
with the first course as its prerequisite --- CNNs and RNNs, BPTT and LSTM, GANs,
sparse coding, dropout and batch normalisation, ResNet and GoogLeNet, CuDNN.}
\edrow{2025}{A single \emph{Deep Learning} course in both the DS and AI clusters. 3+0.}{Four
CLOs in the DS cluster, three in the AI cluster. No outline.}
\verdict{2023's split is genuinely awkward --- the two outlines overlap on
generalization theory, multi-layer perceptrons, backpropagation, CNNs and
autoencoders, so a student taking both repeats a third of the material. 2025
merging them is the right call, but it also means the historical neural-network
material (Hebbian learning, Boltzmann machines, SOMs) quietly disappears.}
\end{coursecard}

\begin{coursecard}{Data Mining}{absent, then a data science pillar}
\edrow{2017}{Does not exist, at either level.}{No course of this name appears in
the 75-course bachelor list or in any MS programme. Association rules and
frequent patterns are taught nowhere.}
\edrow{2023}{BSDS domain core (Semester 5); BSAI domain elective. 3 (2-3).}{Prerequisites
\emph{Advanced Statistics} and \emph{Introduction to Data Science}. Outline:
preprocessing and summary statistics; association rule mining with Apriori and
FP-trees; supervised classification (decision trees, naive Bayes, $k$-NN, SVM);
unsupervised classification ($k$-means, $k$-median, hierarchical and divisive
clustering, Kohonen maps); outlier and anomaly detection; web and social network
mining. Python.}
\edrow{2025}{In both the DS and AI clusters. 3+0.}{Three CLOs, pitched at
knowledge discovery and decision making rather than algorithms.}
\verdict{It appeared from nothing in 2023 and stayed. The overlap with Machine
Learning is severe and is discussed in Part~6 --- roughly half the 2023 Data
Mining outline is also in the 2023 Machine Learning outline.}
\end{coursecard}

\begin{coursecard}{Introduction to Data Science}{the gateway course}
\edrow{2017}{MS only, as \emph{Tools and Techniques in Data Science}, 2+1.}{An
MS(DS) core course. The undergraduate programmes have nothing comparable.}
\edrow{2023}{BSDS domain core, Semester 4. 3 (2-3).}{Prerequisite \emph{Artificial
Intelligence}. Outline: what data science is, datafication, skill sets;
statistical inference, populations and samples, fitting a model, Python;
exploratory data analysis and the data science process; basic algorithms (linear
regression, $k$-NN, $k$-means, naive Bayes); feature generation and selection;
dimensionality reduction with SVD and PCA; mining social network graphs; data
visualization; privacy, security and ethics.}
\edrow{2025}{Data Science cluster, 3+0.}{Three CLOs: understand the principles,
apply EDA and the data science process, apply basic ML algorithms to real
problems of moderate complexity.}
\verdict{Stable in intent across 2023 and 2025 --- it is the process-and-EDA
course, with a taste of modelling. Its prerequisite on \emph{Artificial
Intelligence} in 2023 is odd, since almost nothing in the outline depends on
search or knowledge representation; the real dependencies are programming,
statistics and databases.}
\end{coursecard}

\begin{coursecard}{Natural Language Processing}{the same name, three different courses}
\edrow{2017}{Bachelor elective, 3(3,0), no prerequisite.}{Purely
symbolic and statistical: deterministic and stochastic grammars, parsing
algorithms, CFGs, semantics and semantic roles, corpus methods, $n$-grams and
HMMs, smoothing and backoff, POS tagging and morphology, information retrieval,
the vector space model, precision and recall, information extraction, machine
translation, bag of words. No neural content whatsoever.}
\edrow{2023}{BSAI domain elective; BSDS-adjacent. 3 (2-3).}{Prerequisite
\emph{Artificial Neural Networks \& Deep Learning} --- so the course is now
explicitly downstream of deep learning. The CLOs still read like the 2017 ones.}
\edrow{2025}{In both the DS and AI clusters, 3+0.}{DS-cluster CLOs name
tokenization, POS tagging, parsing, NLTK, sentiment analysis and text
classification. AI-cluster CLOs are broader and add an ethics dimension.}
\verdict{The syllabus of record is still recognisably the 2000s course. The one
real change is the 2023 prerequisite, which correctly places NLP after deep
learning --- but the published content never caught up. Neither 2023 nor 2025
mentions transformers, attention, embeddings or language models anywhere in the
NLP specification.}
\end{coursecard}

\begin{coursecard}{Big Data Analytics}{the most stable course in the cluster}
\edrow{2017}{Bachelor elective, 3(2,1).}{Prerequisites \emph{Probability and
Statistics} and \emph{Programming Fundamentals}. Outline: big data platforms,
Hadoop, storage and analytics, ML algorithms on big data, recommendation,
clustering and classification, graph computing and graph analytics, graphical
models and Bayesian networks, visualization, cognitive mobile analytics.
Leskovec, Rajaraman \& Ullman; Tom White.}
\edrow{2023}{BSDS domain elective. 3 (2-3).}{Prerequisite \emph{Introduction to
Data Science}. Outline moves decisively to Spark: HDFS, MapReduce, YARN, Scala
basics and advanced, RDDs, Spark SQL, ML on Hadoop and Spark, Spark Streaming.}
\edrow{2025}{DS cluster, 3+0.}{Three CLOs, still naming Hadoop, MapReduce and
Apache Spark explicitly --- the only 2025 course in the cluster that names
specific technologies in its outcomes.}
\verdict{Present in all three editions with a coherent story: Hadoop-centric in
2017, Spark-centric from 2023. In 2025 it now shares ground with the new
\emph{Data Engineering} course and the two need deliberate separation.}
\end{coursecard}

\begin{coursecard}{The 2025 arrivals}{Generative AI, Agentic AI, MLOps, Data Engineering}
\edrow{2017}{None of these exist.}{Not as courses, not as topics.}
\edrow{2023}{None of these exist.}{The word ``generative'' does not appear in the
2023 course specifications for this cluster.}
\edrow{2025}{Four new courses, all 3+0.}{\textbf{Generative AI} (in both the DS
and AI clusters): GANs, VAEs and large language models; implementing and
fine-tuning with TensorFlow or PyTorch; ethical and societal implications.
\textbf{Agentic AI} (AI cluster): autonomous agents that reason, plan and use
tools; integration with APIs, databases and external tools; responsible design.
\textbf{MLOps} (AI cluster): end-to-end pipelines for deploying, monitoring and
maintaining models; modern tooling; ethical and security assessment.
\textbf{Data Engineering} (DS cluster): pipelines, ETL, storage systems, data
integration and transformation, architectures for big data and real-time
analytics.}
\verdict{This is the clearest signal of intent in the 2025 booklet. It is also
the hardest part to staff: these four courses have the least settled syllabi, the
fastest-moving tooling, and no published outline to anchor them.}
\end{coursecard}

\subsection{The infrastructure courses}

\begin{coursecard}{Computer Networks}{the most stable course in all three booklets}
\edrow{2017}{Computing core for every degree, 3+1. No prerequisite.}{BSIT
Semester~IV. Outline: protocol architecture, topologies, layered architecture,
physical and data link layer, multiple access, circuit and packet switching, LAN
technologies, wireless, MAC addressing, network devices, IPv4 and IPv6,
addressing, subnetting, CIDR, routing protocols, transport layer, ports and
sockets, flow and congestion control, application layer protocols.}
\edrow{2023}{Computing core for every degree, 3 (2-3). No prerequisite.}{BSIT
Semester~3 --- one semester earlier than 2017. The published outline is
\textbf{word for word identical to the 2017 outline}, down to the phrase ``latest
trends in computer networks''. Only the laboratory allocation changed.}
\edrow{2025}{Major core, Semester~V, 3 credit hours.}{CLOs only. Two of the five
are new in substance: \emph{analyze network security threats and mitigation
techniques} and \emph{design simple network architectures to meet organisational
needs} --- the course is asked to reach past protocol description into design.}
\verdict{Almost nothing. It is compulsory in every degree of every edition, it has
never had a prerequisite, and its 2023 syllabus is a verbatim copy of its 2017
syllabus. If you want a control against which to measure how much the rest of
this document moved, this is it.}
\end{coursecard}

\begin{coursecard}{System \& Network Administration}{compulsory, then compulsory, then optional}
\edrow{2017}{BS-IT core (compulsory), 3 (3,0). Prerequisite \emph{Operating Systems}.}{BSIT
Semester~V. A detailed and frankly practical outline: system administration
components, Microsoft and Linux server environments, server hardware costing,
maintenance contracts and spare parts, data integrity, client--server
configuration, remote console access, comparative analysis of operating systems,
Linux installation and verification, users and groups, software management,
network services and monitoring tools, boot and process management, IP tables and
filtering, securing network traffic, advanced file systems and logs, Bash shell
scripting.}
\edrow{2023}{BSIT domain core, Semester~5, 3 (2-3).}{\textbf{No course contents
are published.} The 2023 booklet specifies only the twenty-nine courses common to
all computing degrees plus the AI, Data Science and Cyber Security clusters ---
there is no Information Technology table at all. A department teaching this course
under HEC 2023 has a title, a credit count and nothing else.}
\edrow{2025}{One elective among fifteen in the Information Technology specialization. 2+1.}{CLOs
only. Notably, two of the five are not technical: \emph{demonstrate
professionalism in managing shared computing resources} and \emph{work
effectively in teams to maintain IT infrastructure}.}
\verdict{It went from a compulsory course with a two-hundred-word syllabus to an
optional course with five outcome statements. For a BSIT department this is the
sharpest loss of specification in the three editions, and the 2017 outline remains
the only usable description of the course that HEC has ever published.}
\end{coursecard}

\begin{coursecard}{Cloud Computing}{the exact opposite trajectory}
\edrow{2017}{Named in the BS-SE elective list at 3-0, and nowhere else.}{It is
\textbf{not} in the seventy-five-course bachelor list and has \textbf{no course
specification} anywhere in the booklet. A BSIT student's nearest equivalent was
\emph{Virtual Systems and Services} (3 cr, prerequisite Programming
Fundamentals) --- a hypervisor course covering VMware, Xen and Hyper-V,
para-virtualised components, and processor support for virtualisation. That is
virtualisation, not cloud: no service models, no elasticity, no billing, no
managed services.}
\edrow{2023}{BSIT domain elective example, 3 (2-3), and in the master course list.}{Also
\textbf{unspecified} --- same gap as System \& Network Administration.}
\edrow{2025}{\textbf{Major core course, Semester~VI, 3 credit hours.} Compulsory
for every BS Computer Science student.}{CLOs: describe IaaS, PaaS and SaaS;
deploy and manage applications in cloud environments; analyze benefits and
challenges; implement basic cloud security and compliance; evaluate cloud
solutions for scalability, cost and performance. It also anchors an entire
specialization --- \emph{Network Infrastructure \& Cloud Computing}, twenty
courses including DevOps, DevSecOps, Infrastructure as Code, Microservices and
Docker, Software Defined Networks, Cloud-Native Application Development, and
Next-Generation Networks and Edge Computing.}
\verdict{From a name with no syllabus to a compulsory course plus a twenty-course
specialization in eight years --- the steepest promotion in the three editions.
And that specialization is where HEC 2025 put \emph{Machine Learning} and
\emph{MLOps}, making it the only cluster outside Data Science and AI that reaches
into this material.}
\end{coursecard}

\begin{notebox}[The specification asymmetry in HEC 2023]
It is worth naming what the three cards above expose. HEC 2023 publishes full
course contents --- introduction, CLOs with Bloom levels, prerequisites, topic
outline and reference books --- for the Artificial Intelligence cluster, the Data
Science cluster and the Cyber Security cluster. It publishes \emph{nothing} for
the Information Technology cluster. Web Technologies, Database Administration \&
Management, System \& Network Administration, Information Technology
Infrastructure, Virtual Systems \& Services, Enterprise Systems, Web Engineering
and Cloud Computing all appear in the BSIT tables and in the master course list,
and none of them has a specification. Computer Networks and Network Security are
specified only because they are common to every computing degree. The practical
consequence is that a BSIT department writing course files for its own domain core
is on its own, and should use the 2017 outlines --- which are detailed, still
broadly accurate, and the only ones that exist.
\end{notebox}

%=====================================================================
\section{What Stays the Same}
%=====================================================================

Underneath the churn there is a stable spine. A department planning a pathway can
rely on all of the following holding across every edition.

\begin{tightnum}
  \item \term{The degree is 130 credit hours} in all three editions, over eight
        semesters, with a final year project.
  \item \term{The programming and data foundation never changes.} Programming
        Fundamentals, Object Oriented Programming, Data Structures and Database
        Systems are compulsory in every edition of every computing degree, and
        they always come in the first three or four semesters.
  \item \term{Artificial Intelligence is the gateway.} Where the cluster is
        taught at all, AI comes first --- as a stated prerequisite in 2023, as
        the earliest compulsory cluster course in 2025.
  \item \term{The cluster is bought with elective credit.} No edition makes
        machine learning compulsory for an IT student. It is always an elective
        choice, which means it is always a departmental decision rather than an
        HEC one.
  \item \term{Three courses appear in all three editions:} Natural Language
        Processing, Big Data Analytics and Computer Vision. If a department has
        taught anything from this cluster since 2017, it is most likely one of
        these.
  \item \term{Computer Networks is the one true constant.} It is compulsory in
        every degree of every edition, it has never carried a prerequisite, and
        its 2023 syllabus is a verbatim reproduction of its 2017 syllabus. Every
        other course in this document moved.
  \item \term{Mathematics support is consistent in intent:} calculus, linear
        algebra and probability and statistics. What changes is whether
        probability is compulsory --- see the warning in Part~9.
  \item \term{Bloom's taxonomy labelling of course learning outcomes} is used in
        all three, so outcomes remain comparable even when outlines are not.
\end{tightnum}

%=====================================================================
\section{What Changes, and Why It Matters}
%=====================================================================

\subsection{2017 to 2023: the cluster comes down to the bachelor level}

\begin{tight}
  \item \textbf{Artificial Intelligence becomes compulsory for every computing
        degree}, BSIT included, in Semester 3. This is the change that makes an
        undergraduate data science pathway possible at all.
  \item \textbf{Eight courses appear from nothing}: Machine Learning, Data
        Mining, Introduction to Data Science, Data Visualization, Data
        Warehousing \& Business Intelligence, Reinforcement Learning, Knowledge
        Representation \& Reasoning and Programming for AI all become
        undergraduate courses with published outlines.
  \item \textbf{Two new degrees}, BSAI and BSDS, give the cluster a home and,
        more usefully for everyone else, give it a specified curriculum that
        other degrees can borrow from.
  \item \textbf{Prerequisite chains are published and enforced.} 2023 is the only
        edition that states, course by course, what must come first. Those chains
        are the backbone of the pathway in Part~7.
  \item \textbf{Laboratory hours are made systematic}: nearly every cluster
        course becomes 3 (2-3) --- two hours of lecture and three of laboratory.
  \item \textbf{But the Information Technology cluster is left unspecified.}
        2023 writes full course contents for the AI, Data Science and Cyber
        Security clusters and none at all for the IT one --- System \& Network
        Administration, Cloud Computing, IT Infrastructure, Web Technologies,
        Database Administration and the rest have a title and a credit count and
        nothing more.
\end{tight}

\subsection{2023 to 2025: from degrees to specializations}

\begin{tight}
  \item \textbf{The separate degrees go away} --- at least in this booklet, which
        covers Computer Science only. Data Science and Artificial Intelligence
        become two of fourteen specializations, each defined by a cluster of
        15 elective courses from which a student takes seven or eight.
  \item \textbf{The specialization may be named on the degree or the transcript},
        or not at all: a university may run a generic BSCS in which students take
        electives from any cluster.
  \item \textbf{Seven genuinely new courses} enter the cluster: Generative AI,
        Agentic AI, MLOps, Data Engineering, Information Retrieval, Data Ethics \&
        Security and Optimization Techniques.
  \item \textbf{Prerequisites disappear.} The 2025 booklet publishes no
        prerequisite for any specialization course. The elective pool is formally
        flat, and sequencing becomes entirely a departmental responsibility.
  \item \textbf{Topic outlines disappear too.} Only CLOs remain. This buys
        flexibility at the cost of comparability between universities.
  \item \textbf{Laboratory hours mostly disappear from the cluster}: the Data
        Science and AI specialization courses are listed as 3+0, where 2023 gave
        them three laboratory hours each.
  \item \textbf{Two new compulsory blocks appear}: a professional certification
        requirement worth three credit hours, and a three-credit field experience
        or internship that explicitly cannot be substituted with coursework.
  \item \textbf{Cloud Computing becomes compulsory} for every BS Computer Science
        student, in Semester~VI, having had no specification at all in 2017. In
        the same move \emph{System \& Network Administration} and \emph{IT
        Infrastructure} --- compulsory for BSIT in both earlier editions --- drop
        to optional.
  \item \textbf{Infrastructure gets its own specialization, and it reaches into
        this cluster.} \emph{Network Infrastructure \& Cloud Computing} lists
        twenty courses including DevOps, DevSecOps, Infrastructure as Code,
        Microservices and Docker, Software Defined Networks and Edge Computing
        --- and also \emph{Machine Learning} and \emph{MLOps}.
\end{tight}

\begin{notebox}[Naming drift worth tracking in your own scheme of studies]
\begin{tight}
  \item \emph{Advanced Statistics} (2023) $\rightarrow$ \emph{Computational
        Statistics} (2025)
  \item \emph{Artificial Neural Networks \& Deep Learning} (2023) $\rightarrow$
        \emph{Deep Learning} (2025)
  \item \emph{Data Warehousing \& Business Intelligence} (2023) $\rightarrow$
        \emph{Business Intelligence} + \emph{Data Engineering} (2025)
  \item \emph{Tools and Techniques in Data Science} --- an MS core course in
        2017, an undergraduate DS elective in 2025, absent in 2023
  \item \emph{Digital Image Processing} (2017, 2023) $\rightarrow$ absorbed into
        \emph{Computer Vision} (2025)
  \item \emph{Text Mining} and \emph{Topics in Data Science} (2023) $\rightarrow$
        dropped; nearest 2025 equivalents are \emph{Information Retrieval} and
        \emph{NLP}
  \item \emph{Virtual Systems and Services} (2017, 2023) $\rightarrow$ dropped;
        its ground is taken by \emph{Cloud Computing} and, in the 2025
        infrastructure cluster, by \emph{Microservices Architecture and Docker
        Containers}
  \item \emph{IT Infrastructure} (2017) $\rightarrow$ \emph{Information
        Technology Infrastructure} (2023, 2025) --- same course, longer name
  \item \emph{Cloud Computing} (2025 major core) and \emph{Cloud Infrastructure
        and Services} (2025 IT cluster) are \textbf{two different courses} in the
        same booklet with heavily overlapping outcomes; see Part~6
\end{tight}
\end{notebox}

%=====================================================================
\section{Where the Courses Overlap}
%=====================================================================

The cluster contains real duplication. A department that offers several of these
courses without deciding the boundaries will teach decision trees four times and
transformers never. The table gives the overlap and the split that resolves it.

\begin{longtable}{@{}L{4.4cm}L{5.2cm}L{6.0cm}@{}}
\toprule
\rowcolor{primary!15}
\textbf{Pair} & \textbf{What both claim} & \textbf{Recommended split} \\
\midrule
\endhead
Artificial Intelligence \newline vs Machine Learning &
Supervised, unsupervised and reinforcement learning are named in the AI outlines
of both 2017 and 2023, and ANNs are named in the 2023 AI outline &
Keep AI symbolic: search, constraint satisfaction, adversarial search, knowledge
representation, logic, planning, uncertainty. Give ML everything statistical, and
let AI cover learning in one week as a signpost, not a unit \\
\rowcolor{lightbg}
Machine Learning \newline vs Data Mining &
Decision trees, naive Bayes, $k$-NN, SVM, $k$-means and hierarchical clustering
appear in both 2023 outlines --- roughly half of Data Mining &
ML owns prediction: generalization, model selection, evaluation, ensembles,
regularisation. Data Mining owns discovery: preprocessing at scale, association
rules and frequent patterns, anomaly and outlier detection, graph and social
network mining \\
Introduction to Data Science \newline vs Machine Learning &
Linear regression, $k$-NN, $k$-means and naive Bayes are in the 2023 Intro to DS
outline &
Intro to DS owns the process: framing, data acquisition and cleaning, EDA,
feature engineering, dimensionality reduction, communication, ethics. It should
demonstrate two or three models, never teach the model zoo \\
\rowcolor{lightbg}
ANN \& Deep Learning \newline vs Deep Learning (2023 only) &
Generalization theory, multi-layer perceptrons, backpropagation, CNNs and
autoencoders are in both outlines &
Do not offer both. Offer the 2025 merged \emph{Deep Learning} instead \\
Big Data Analytics \newline vs Data Engineering (2025) &
Storage, distributed processing and pipelines &
Big Data Analytics owns the distributed compute model: HDFS, MapReduce, Spark,
streaming. Data Engineering owns the pipeline: ingestion, ETL, orchestration,
schema and contract design, warehousing \\
\rowcolor{lightbg}
NLP vs Text Mining \newline vs Information Retrieval &
Vector space model, text classification, precision and recall appear in all three &
Offer NLP as the language course (representation, sequence models, generation);
offer Information Retrieval only if search ranking and indexing are taught
properly. Do not offer Text Mining as well \\
Data Visualization \newline vs Business Intelligence &
Charts, dashboards, decision support &
Visualization owns perception, chart selection and the grammar of graphics; BI
owns the enterprise stack: dimensional modelling, OLAP, reporting, governance \\
\rowcolor{lightbg}
Generative AI vs Deep Learning &
GANs and VAEs are named in both 2025 CLO sets &
Deep Learning owns architectures and training; Generative AI owns large
pretrained models, prompting, fine-tuning, retrieval augmentation and evaluation
of generated output \\
Cloud Computing \newline vs Cloud Infrastructure and Services (2025) &
Service models, deployment, management, compliance --- these are two courses in
the same booklet, one a major core and one an IT elective, with four of five
outcomes in common &
Do not offer both to the same student. If the IT specialization is being run,
treat \emph{Cloud Infrastructure and Services} as satisfied by the core
\emph{Cloud Computing} and spend the slot elsewhere \\
\rowcolor{lightbg}
Cloud Computing \newline vs Virtual Systems \& Services (2017, 2023) &
Hypervisors, virtual machines, resource pooling &
Virtualisation is one chapter of cloud computing, not a peer of it. Where both
exist in a scheme of studies, retire \emph{Virtual Systems \& Services} \\
System \& Network Administration \newline vs IT Infrastructure &
Servers, storage, operating systems, availability, monitoring --- the 2017
outlines overlap on most of the middle of both courses &
Administration owns the hands-on: installation, users and groups, services,
firewall rules, shell scripting, logs. Infrastructure owns the architectural:
availability and performance patterns, data centres, capacity, service delivery
and support processes, lifecycle \\
\rowcolor{lightbg}
Cloud Computing \newline vs Big Data Analytics and MLOps &
Elastic compute, cluster provisioning, managed services &
Cloud owns the platform and its economics: service models, provisioning, cost,
security and compliance. Big Data Analytics owns the distributed compute model
that runs on it; MLOps owns the model lifecycle that runs on both \\
\bottomrule
\end{longtable}

\begin{recbox}[The one sequencing rule that follows from this table]
If a scheme of studies contains Cloud Computing \emph{and} Big Data Analytics
\emph{and} MLOps, teach them in that order. Each one assumes the previous one's
platform, and taught in the wrong order the later course spends its first three
weeks re-teaching the earlier one. This is the only infrastructure sequencing
constraint in the whole document, and no HEC edition states it.
\end{recbox}

% Keep the Part 7 heading, its opening paragraph and the pathway figure on one
% page: the figure is the argument, and a heading stranded on the page before
% it makes the reader turn back and forth.
\Needspace*{0.80\textheight}
%=====================================================================
\section{The Pathway}
%=====================================================================

The diagram reconciles the three editions into one dependency graph. Solid arrows
are prerequisites published in the 2023 booklet, which is the only edition that
states them. Dashed arrows are dependencies that are real but unpublished --- they
are what a department must enforce itself, particularly under 2025, where no
prerequisites are given at all.

% Forced in place: the diagram is the argument of this section and must not
% drift to a float page away from the text that reads it.
\begin{figure}[H]
\centering
\resizebox{\ifdim\width>\linewidth \linewidth\else \width\fi}{!}{%
\begin{tikzpicture}[
  font=\footnotesize,
  box/.style={rounded corners=3pt, draw=#1, fill=#1!8, text=#1!45!black,
              text width=3.15cm, align=center, minimum height=1.05cm,
              font=\footnotesize, inner sep=3pt},
  tier/.style={font=\footnotesize\bfseries, text=primary, anchor=east,
               align=right},
  band/.style={fill=#1, rounded corners=3pt, draw=none},
  pre/.style={-{Stealth[length=2.4mm]}, line width=0.9pt, primary!70},
  soft/.style={-{Stealth[length=2.4mm]}, line width=0.9pt, neutral,
               dash pattern=on 2.2pt off 1.8pt}]

% ---- tier bands, drawn first so every arrow and box sits on top
\fill[band=primary!4]   (-0.30,-0.75) rectangle (16.9,0.75);
\fill[band=secondary!5] (-0.30,-3.15) rectangle (16.9,-1.65);
\fill[band=goldhl!7]    (-0.30,-5.55) rectangle (16.9,-4.05);
\fill[band=greenhl!6]   (-0.30,-7.95) rectangle (16.9,-6.45);
\fill[band=purplehl!6]  (-0.30,-10.35) rectangle (16.9,-8.85);

\node[tier] at (-0.50,0)      {Foundation\\\mdseries Sem 1--3};
\node[tier] at (-0.50,-2.4)   {Gateway\\\mdseries Sem 3--5};
\node[tier] at (-0.50,-4.8)   {Entry\\\mdseries Sem 5};
\node[tier] at (-0.50,-7.2)   {Core methods\\\mdseries Sem 6};
\node[tier] at (-0.50,-9.6)   {Depth\\\mdseries Sem 6--7};

% =====================  Tier 0 -- foundation  =====================
\node[box=neutral] (pf)  at (2.4,0)  {Programming\\Fundamentals};
\node[box=neutral] (dsb) at (6.6,0)  {Data Structures\\+ Database Systems};
\node[box=neutral] (la)  at (10.8,0) {Linear Algebra\\+ Calculus};
\node[box=neutral] (ps)  at (15.0,0) {Probability\\\& Statistics};

% =====================  Tier 1 -- gateway  ========================
\node[box=primary] (ai) at (4.5,-2.4)  {\textbf{Artificial}\\\textbf{Intelligence}};
\node[box=tealhl]  (cn) at (11.5,-2.4) {\textbf{Computer Networks}};

% =====================  Tier 2 -- entry  ==========================
\node[box=secondary] (ml)  at (2.4,-4.8)  {\textbf{Machine}\\\textbf{Learning}};
\node[box=secondary] (ids) at (6.6,-4.8)  {\textbf{Introduction to}\\\textbf{Data Science}};
\node[box=tealhl]    (sna) at (10.8,-4.8) {System \& Network\\Administration};
\node[box=tealhl]    (cc)  at (15.0,-4.8) {\textbf{Cloud Computing}};

% =====================  Tier 3 -- core methods  ===================
\node[box=goldhl] (dl)  at (2.4,-7.2)  {\textbf{Deep Learning}};
\node[box=goldhl] (dm)  at (7.0,-7.2)  {Data Mining};
\node[box=goldhl] (bda) at (12.0,-7.2) {Big Data Analytics\\/ Data Engineering};

% =====================  Tier 4 -- depth  ==========================
\node[box=purplehl] (nlp) at (2.4,-9.6)  {NLP};
\node[box=purplehl] (cv)  at (6.6,-9.6)  {Computer Vision};
\node[box=purplehl] (gai) at (10.8,-9.6) {Generative AI\\$\rightarrow$ Agentic AI};
\node[box=purplehl] (mlo) at (15.0,-9.6) {\textbf{MLOps}};

% ---- prerequisites published in HEC 2023 (solid)
\draw[pre] (pf.south)       -- (ai.north west);
\draw[pre] (dsb.south)      -- (ai.north east);
\draw[pre] (ai.south west)  -- (ml.north);
\draw[pre] (ai.south east)  -- (ids.north west);
\draw[pre] (ml.south)       -- (dl.north);
\draw[pre] (ids.south)      -- (dm.north west);
\draw[pre] (ids.south east) -- (bda.north west);
\draw[pre] (dl.south)       -- (nlp.north);

% ---- real but unpublished dependencies (dashed). Only edges that can be
% drawn without crossing another edge or a box are drawn; the rest are
% stated in the note beneath, which is more honest than a tangle.
\draw[soft] (cn.south)       -- (sna.north);
\draw[soft] (cn.south east)  -- (cc.north west);
\draw[soft] (cc.south)       -- (bda.north east);
\draw[soft] (ml.south east)  -- (dm.north west);
\draw[soft] (dl.south east)  -- (cv.north west);
\draw[soft] (dl.south east)  -- (gai.north west);
\draw[soft] (bda.south east) -- (mlo.north west);

% ---- legend
\draw[pre]  (0.1,-11.0) -- (1.1,-11.0);
\node[anchor=west, font=\scriptsize, text=neutral] at (1.2,-11.0)
     {prerequisite published in HEC 2023};
\draw[soft] (7.6,-11.0) -- (8.6,-11.0);
\node[anchor=west, font=\scriptsize, text=neutral] at (8.7,-11.0)
     {real dependency, published in no edition};
\node[box=tealhl, minimum height=0.42cm, text width=2.05cm, font=\scriptsize]
     at (13.9,-11.0) {infrastructure};
\node[anchor=north west, align=left, text width=16.8cm, font=\scriptsize,
      text=neutral] at (0.1,-11.35)
     {\textbf{The foundation band underpins everything below it} and its arrows
      are omitted for legibility: linear algebra and calculus are assumed by
      Machine Learning and Deep Learning, and probability and statistics by every
      course from the Entry band down. \textbf{The teal lane is the
      infrastructure the rest of the graph runs on}: it joins the data lane at Big
      Data Analytics and again at MLOps, which also assumes Machine Learning.
      \textbf{HEC 2025 publishes no prerequisite for any specialization course},
      so under that edition every arrow here is a departmental responsibility
      rather than a curricular rule. Computational Statistics, Data
      Visualization, Business Intelligence, Reinforcement Learning, Information
      Retrieval, Speech Processing and Data Ethics \& Security are satellites of
      this graph: each hangs off one course above it and none is a prerequisite
      for anything else.};
\end{tikzpicture}}
\caption{The AI, data science and infrastructure pathway, reconciled across the
three HEC editions. The five bands are the order in which the material can
actually be learned; the teal lane is the platform the data lane runs on.
Semester ranges assume the HEC 2023 BSIT structure, in which Artificial
Intelligence and Computer Networks are compulsory by Semester~3 and elective
slots open in Semester~5.}\label{fig:pathway}
\end{figure}

\subsection{Reading the pathway}

\begin{tight}
  \item \textbf{Nothing in the cluster can start before Semester 5} for a BSIT
        student, because the earliest elective slots are in Semester 5 and
        \emph{Artificial Intelligence} --- the published prerequisite for both
        Machine Learning and Introduction to Data Science --- is not finished
        until the end of Semester 3.
  \item \textbf{The chain has depth four}: AI $\rightarrow$ Machine Learning
        $\rightarrow$ Deep Learning $\rightarrow$ NLP. With elective slots only
        in Semesters 5, 6 and 7 this exactly fills the available room, and it is
        why NLP cannot sensibly be offered before the final year.
  \item \textbf{Two entry points, not one.} A student may enter through
        \emph{Machine Learning} (the modelling route) or through
        \emph{Introduction to Data Science} (the process route). The second is
        the better entry for IT students, because it starts from data handling
        and databases, which they already have.
  \item \textbf{The frontier tier is optional and should be treated as such.}
        Generative AI, Agentic AI and MLOps are genuinely valuable and genuinely
        hard to staff. A pathway that delivers the first three tiers well is
        better than one that offers all five badly.
\end{tight}

%=====================================================================
\section{Semester-Wise Schemes}
%=====================================================================

Four schemes follow. All of them leave the compulsory parts of the degree
untouched and spend only elective credit, so none of them requires HEC approval
of a new programme --- they are choices a department makes within an approved
curriculum.

\subsection{Scheme A --- BSIT under HEC 2023, full data science track}

Seven domain electives, twenty-one credit hours, all spent on the cluster. This
produces a genuine concentration and satisfies every prerequisite the 2023
booklet publishes. It also spends the entire elective budget, which is the
trade-off discussed after the table.

\begin{longtable}{@{}L{1.5cm}L{6.6cm}L{2.2cm}L{5.0cm}@{}}
\toprule
\rowcolor{primary!15}
\textbf{Sem} & \textbf{Course} & \textbf{Credits} & \textbf{Status and prerequisite} \\
\midrule
\endhead
1--2 & Programming Fundamentals; Object Oriented Programming; Database Systems;
Multivariable Calculus; Linear Algebra & as per BSIT &
Already compulsory. No change \\
\rowcolor{lightbg}
3 & Data Structures; \textbf{Artificial Intelligence}; Probability \& Statistics &
as per BSIT & Already compulsory. \emph{Artificial Intelligence} is the gateway,
and it is already in the core \\
4 & --- & --- & No elective slots. Domain cores only \\
\rowcolor{lightbg}
5 & \textbf{Introduction to Data Science} \newline (Domain Elective 1) & 3 (2-3) &
Prerequisite: Artificial Intelligence \checkmark\ Sem 3 \\
5 & \textbf{Machine Learning} \newline (Domain Elective 2) & 3 (2-3) &
Prerequisite: Artificial Intelligence \checkmark\ Sem 3 \\
\rowcolor{lightbg}
6 & \textbf{Data Mining} (Domain Elective 3) & 3 (2-3) &
Prerequisite: Introduction to Data Science \checkmark\ Sem 5 \\
6 & \textbf{Deep Learning} (Domain Elective 4) & 3 (2-3) &
Prerequisite: Machine Learning \checkmark\ Sem 5 \\
\rowcolor{lightbg}
6 & \textbf{Big Data Analytics} (Domain Elective 5) & 3 (2-3) &
Prerequisite: Introduction to Data Science \checkmark\ Sem 5 \\
6 & \textbf{Data Visualization} (Domain Elective 6) & 3 (2-3) &
Pairs with Data Mining; no hard prerequisite \\
\rowcolor{lightbg}
7 & \textbf{Natural Language Processing} \newline (Domain Elective 7) & 3 (2-3) &
Prerequisite: Deep Learning \checkmark\ Sem 6. Substitute Generative AI or MLOps
if staffing favours it \\
7--8 & \textbf{Final Year Project I and II} & 2 + 4 &
Already compulsory. Supervise within the strand \\
\bottomrule
\end{longtable}

\begin{warnbox}[What Scheme A costs]
Spending all seven domain electives on the cluster means a BSIT graduate takes
\emph{no} Cloud Computing, Network Security, Web Engineering, Enterprise Systems
or Virtual Systems \& Services --- the courses the 2023 booklet itself names as
BSIT elective examples. \textbf{Losing Cloud Computing is the one that hurts},
because Big Data Analytics in Semester~6 then has to teach its own platform from
scratch, and because HEC 2025 has since made Cloud Computing compulsory for every
computing student. Computer Networks (Semester~3), System \& Network
Administration and IT Infrastructure survive --- they are BSIT domain \emph{core},
not electives --- so the graduate is not left without infrastructure entirely.
That may be the right call for a department building a data specialization, but it
should be a decision, not an accident. Scheme~B below is the balanced alternative.
\end{warnbox}

\subsection{Scheme B --- BSIT under HEC 2023, balanced track (recommended)}

Four of the seven electives on the cluster, three left for IT depth. This is the
scheme to adopt if the department wants every BSIT graduate to be competent with
data without turning the degree into a data science degree.

\begin{longtable}{@{}L{1.5cm}L{6.6cm}L{2.2cm}L{5.0cm}@{}}
\toprule
\rowcolor{primary!15}
\textbf{Sem} & \textbf{Course} & \textbf{Credits} & \textbf{Note} \\
\midrule
\endhead
3 & \textbf{Artificial Intelligence} (core); \textbf{Computer Networks} (core) &
3 (2-3) each & Both compulsory already. The two gateways, one per lane \\
\rowcolor{lightbg}
5 & \textbf{Introduction to Data Science} (DE 1) & 3 (2-3) &
The entry point. Process, EDA, feature work, ethics \\
5 & \textbf{Cloud Computing} (DE 2) & 3 (2-3) &
The platform course. Put it here, not later: Big Data Analytics and MLOps both
assume it, and HEC 2025 has since made it compulsory for every computing student \\
\rowcolor{lightbg}
5 & \emph{(System \& Network Administration is a domain core in Semester~5 --- it
costs no elective credit)} & 3 (2-3) &
Follows Computer Networks. Already compulsory for BSIT \\
6 & \textbf{Machine Learning} (DE 3) & 3 (2-3) &
Prerequisite: Artificial Intelligence \checkmark\ Sem 3 \\
\rowcolor{lightbg}
6 & \textbf{Big Data Analytics} (DE 4) & 3 (2-3) &
Prerequisite: Introduction to Data Science \checkmark\ Sem 5. Now has its platform
from Cloud Computing \checkmark\ Sem 5 \\
6 & IT elective --- e.g.\ Network Security (DE 5) & 3 (2-3) &
Keeps the IT identity of the degree \\
\rowcolor{lightbg}
6 & IT elective --- e.g.\ Web Engineering (DE 6) & 3 (2-3) & \\
7 & \textbf{Deep Learning} (DE 7) & 3 (2-3) &
Prerequisite: Machine Learning \checkmark\ Sem 6. Substitute \textbf{Data Mining}
if the departmental strength is statistics rather than neural networks \\
\rowcolor{lightbg}
7--8 & Final Year Project I and II & 2 + 4 &
At least one project topic per supervisor drawn from the strand \\
\bottomrule
\end{longtable}

\begin{recbox}[Why this is the default recommendation]
Four cluster courses --- Introduction to Data Science, Machine Learning, Big Data
Analytics and Deep Learning --- is the smallest set that produces a graduate who
can frame a data problem, build and evaluate a model, run it at scale and
understand what a neural network is doing. The fifth slot goes to \textbf{Cloud
Computing}, which is not a concession to the IT side of the degree but a
prerequisite for the data side: it is what Big Data Analytics runs on, and it is
the course HEC 2025 later made compulsory for everyone. Two IT electives still
survive on top of a domain core that already contains Computer Networks, System \&
Network Administration and IT Infrastructure, so the degree still reads clearly as
Information Technology.
\end{recbox}

\subsection{Scheme C --- BSCS with an IT specialization under HEC 2025}

This is the harder case. The 2025 booklet's \emph{Information Technology}
specialization cluster contains fifteen courses --- Web Technologies, Cyber
Security, Database Administration, System \& Network Administration, IT
Infrastructure, Network Security, Data Communication, Wireless \& Mobile
Networks, Cloud Infrastructure, Email Systems, Ethical Hacking, Security Policies,
IT Project Management, Systems Analysis \& Design and Technology Lifecycle
Planning --- and \emph{not one of them is an AI or data course}. Taking the named
IT specialization strictly means taking none of this cluster.

\begin{finding}[Three ways out, and the booklet supplies all of them]
First, universities may run a \textbf{generic BS Computer Science} in which
students take electives from any cluster, with no specialization named on the
degree or transcript. Second, HEIs \textbf{may add courses to a cluster} after
approval of their own statutory bodies --- this route keeps the specialization
title. Third, and most simply, \textbf{specialization~12, Network Infrastructure
\& Cloud Computing, already contains Machine Learning and MLOps} alongside DevOps,
DevSecOps, Infrastructure as Code, Software Defined Networks, Microservices and
Docker, and Edge Computing. For a student who wants infrastructure with real data
content, that named specialization is a better fit than the Information Technology
one and needs no local approval at all.
\end{finding}

\begin{notebox}[Two things HEC 2025 gives this student for free]
\textbf{Computer Networks} is a major core course in Semester~V and \textbf{Cloud
Computing} is a major core course in Semester~VI. Neither costs an elective slot.
That is a real improvement on 2023 for anyone building a data track: the platform
half of the pathway arrives compulsorily, and all eight elective slots stay
available. It also means \emph{Cloud Infrastructure and Services} in the IT
cluster is largely redundant for these students --- see Part~6.
\end{notebox}

\begin{longtable}{@{}L{1.7cm}L{6.4cm}L{2.2cm}L{5.0cm}@{}}
\toprule
\rowcolor{primary!15}
\textbf{Sem} & \textbf{Course} & \textbf{Credits} & \textbf{Note} \\
\midrule
\endhead
I--IV & Programming Fundamentals; OOP; Digital Logic Design; Data Structures;
Database Systems; Operating Systems; Software Engineering; Computer Organization;
Design \& Analysis of Algorithms & as per BSCS &
Compulsory. Calculus is IDS-I, Linear Algebra is IDS-II \\
\rowcolor{lightbg}
V & \textbf{Artificial Intelligence}; \textbf{Computer Networks} (Major core) & 3 each &
Both compulsory, and both free of elective credit. The two gateways, but two
years later than under 2023 \\
V & \textbf{Probability \& Statistics} as \textbf{IDS-III} & 3 &
\textbf{Must be chosen deliberately} --- see the warning below \\
\rowcolor{lightbg}
VI & \textbf{Cloud Computing} (Major core) & 3 &
Compulsory. Costs no elective slot, and arrives exactly when the data electives
start needing a platform \\
VI & \textbf{Introduction to Data Science} (Elective I) & 3 & DS cluster \\
\rowcolor{lightbg}
VI & \textbf{Machine Learning} (Elective II) & 3 &
DS, AI \emph{and} Network Infrastructure clusters \\
VI & System \& Network Administration (Elective III) & 2+1 & IT cluster \\
\rowcolor{lightbg}
VI & Web Technologies (Elective IV) & 2+1 & IT cluster \\
VII & \textbf{Deep Learning} (Elective V) & 3 & DS and AI clusters \\
\rowcolor{lightbg}
VII & \textbf{Big Data Analytics} (Elective VI) & 3 &
Its platform is already in place from Semester~VI \\
VII & \textbf{MLOps} (Elective VII) & 3 &
Where the two lanes meet. In the AI \emph{and} Network Infrastructure clusters,
so it counts either way \\
\rowcolor{lightbg}
VII & IT Infrastructure (Elective VIII) & 2+1 &
IT cluster. Do \textbf{not} also take \emph{Cloud Infrastructure and Services} ---
the core Cloud Computing course already covers it \\
VII & \textbf{Professional Certification} & 3 &
Compulsory in 2025. Align it with the strand --- the booklet's own AI
specialization annexe recommends MOOC routes for exactly these courses \\
VIII & \textbf{Final Year Project} + \textbf{Field Experience} & 6 + 3 &
Both compulsory. The internship cannot be substituted with coursework \\
\bottomrule
\end{longtable}

\begin{warnbox}[The probability trap in HEC 2025]
Under HEC 2023, Probability \& Statistics is a required mathematics course in
Semester 3 for both BSIT and BSCS. Under HEC 2025 it is \textbf{not compulsory}:
only Calculus and Linear Algebra are mandatory interdisciplinary courses.
Probability \& Statistics is item four in an optional pool of eighteen, and the
two free interdisciplinary slots (IDS-III and IDS-IV) both fall in
\textbf{Semester V}. A student who does not choose it will arrive at Machine
Learning in Semester VI with no statistics at all. Any department running a data
science or AI specialization under the 2025 curriculum should make Probability \&
Statistics the mandated IDS-III for those students, and should consider petitioning
to move it earlier.
\end{warnbox}

\subsection{Scheme D --- the minimum viable strand}

For a department that cannot yet staff a full track, three courses in this order
still produce a usable graduate. Everything else can wait.

\begin{longtable}{@{}L{1.2cm}L{4.9cm}L{9.6cm}@{}}
\toprule
\rowcolor{primary!15}
\textbf{\#} & \textbf{Course} & \textbf{What the graduate can do that they could not before} \\
\midrule
\endhead
1 & Introduction to Data Science &
Frame a question as a data problem, audit and clean a dataset, explore it
honestly, and know when the answer is not in the data \\
\rowcolor{lightbg}
2 & Machine Learning &
Build a predictive model, choose the right metric, avoid leakage, and read a
result critically \\
3 & One applied course --- Big Data Analytics, Data Mining or Deep Learning &
Work at the scale, or with the model family, that the local job market actually
asks for \\
\bottomrule
\end{longtable}

%=====================================================================
\section{Risks and Implementation Notes}
%=====================================================================

\subsection{Things that will go wrong if they are not planned}

\begin{tightnum}
  \item \textbf{Prerequisites will not enforce themselves under HEC 2025.} The
        booklet publishes none. Adopt the 2023 chains explicitly in your own
        scheme of studies and encode them in the student information system.
  \item \textbf{Laboratory hours will quietly disappear.} 2023 gave every cluster
        course three laboratory hours; 2025 lists them as 3+0. This cluster
        cannot be taught from a whiteboard. If you adopt the 2025 structure,
        write the laboratory back in as 2+1.
  \item \textbf{Machine Learning and Data Mining will be taught twice.} Assign
        the boundary in Part~6 to named course owners before the first offering,
        not after the first student complaint.
  \item \textbf{The NLP course will be a decade out of date.} No edition's
        published NLP specification mentions embeddings, attention, transformers
        or language models. The course code and title are fine; the content must
        be written locally.
  \item \textbf{Statistics will be missing.} See the warning above for HEC 2025.
        Even under 2023, one Probability \& Statistics course in Semester~3 is
        thin support for Semester~6 machine learning; a Computational Statistics
        elective is worth the slot if one is free.
  \item \textbf{The frontier courses will be assigned to whoever is available.}
        Generative AI, Agentic AI and MLOps have no published outline in any
        edition. Do not offer one without a named owner who has built something
        in it.
  \item \textbf{The IT courses have no syllabus to inherit.} HEC 2023 specifies
        the AI, Data Science and Cyber Security clusters and leaves the
        Information Technology one completely unspecified; HEC 2025 gives every
        specialization course five outcome statements and nothing else. For
        System \& Network Administration, IT Infrastructure and Virtual Systems \&
        Services, the 2017 outlines are the only detailed descriptions HEC has
        ever published, and they are still broadly sound. Use them rather than
        starting from a blank page.
  \item \textbf{Cloud Computing will be scheduled too late.} It is the platform
        for Big Data Analytics, Data Engineering and MLOps, so a scheme that puts
        it in the final year forces each of those courses to teach its own
        infrastructure from scratch. Put it no later than the semester before the
        first of them.
  \item \textbf{Cloud will be taught twice under HEC 2025.} \emph{Cloud
        Computing} is a compulsory major course and \emph{Cloud Infrastructure and
        Services} is an IT specialization elective, and four of their five
        outcomes are the same. Bar the combination in the scheme of studies.
\end{tightnum}

\subsection{What to do with the 2025 certification requirement}

The 2025 booklet makes three credit hours of professional certification a
condition of degree completion, and its own annexe recommends MOOC routes aligned
to the AI specialization --- a foundational block (probability, discrete
mathematics, data structures) and a core AI and machine learning block mapped
course by course to Coursera and edX offerings. For a student on this pathway the
certification block is best spent reinforcing the strand rather than diversifying
away from it: a machine learning or cloud data engineering certification
completed alongside Semester~VII electives both satisfies the requirement and
strengthens the final year project.

\subsection{A note on the three editions coexisting}

Universities do not switch curricula cleanly. In practice a department may be
running 2023 for its current cohorts, planning 2025 for the next intake, and
still carrying 2017-era course files in its quality-assurance records. When that
happens the mapping that matters is the one in Part~2: it lets a course approved
under one edition be recognised under another. Two rules make the mapping safe.
Match on \emph{outline}, not on title --- \emph{Advanced Statistics} and
\emph{Computational Statistics} are the same course, while the 2017 and 2025
courses both called \emph{Natural Language Processing} are not. And where an
edition splits a course, treat the pair as one --- a student who took 2023's
\emph{ANN \& Deep Learning} has met 2025's \emph{Deep Learning} requirement.

%=====================================================================
\section{The Published Course Contents}
%=====================================================================

Everything below is what the three booklets actually say, reproduced rather than
summarised, so that a department writing a course file can lift the wording
directly and see at a glance which edition gives it most to work with.

Three things to know before reading. \textbf{HEC 2017} publishes credit hours, a
prerequisite, CLOs with Bloom levels, a topic outline and reference books for all
seventy-five bachelor courses it lists. \textbf{HEC 2023} publishes the same, plus
a course-introduction paragraph, but \emph{only for forty-six courses} --- those
common to every computing degree, plus the Artificial Intelligence, Data Science
and Cyber Security cluster tables. \textbf{HEC 2025} publishes course learning
outcomes and nothing else: no prerequisite, no outline, no reading list, for any
of its two hundred and fifteen specialization courses.

Where an edition says nothing, this is recorded rather than passed over --- the
gaps are as informative as the text, and several of them are the reason a
department has to write the course itself.

\subsection{The AI and data science courses}

\begin{coursecard}{Artificial Intelligence}{the gateway course}
\edrow{2017}{3+1 \;$\bullet$\; prerequisite: Data Structures and Algorithms}{Introduction (Introduction, basic component of AI, Identifying AI systems, branches of AI, etc.); Reasoning and Knowledge Representation (Introduction to Reasoning and Knowledge Representation, Propositional Logic, First order Logic); Problem Solving by Searching (Informed searching, Uninformed searching, Local searching.); Constraint Satisfaction Problems; Adversarial Search (Min-max algorithm, Alpha beta pruning, Game-playing); Learning (Unsupervised learning, Supervised learning, Reinforcement learning) ;Uncertainty handling (Uncertainty in AI, Fuzzy logic); Recent trends in AI and applications of AI algorithms (trends, Case study of AI systems, Analysis of AI systems) \par\vspace{2pt}{\footnotesize\itshape Set texts: Stuart Russell and Peter Norvig, Artificial Intelligence. A Modern Approach, 3rd edition, Prentice Hall, Inc., 2010; Hart, P.E., Stork, D.G. and Duda, R.O., 2001. Pattern classification. John Willey \& Sons; Luger, G.F. and Stubblefield, W.A., 2009. AI algorithms, data structures, and idioms in Prolog, Lisp, and Java. Pearson Addison-Wesley.}}
\edrow{2023}{3 (2-3) \;$\bullet$\; prerequisite: Object Oriented Programming}{An Introduction to Artificial Intelligence and its applications towards Knowledge Based Systems; Introduction to Reasoning and Knowledge Representation, Problem Solving by Searching (Informed searching, Uninformed searching, Heuristics, Local searching, Min- max algorithm, Alpha beta pruning, Game-playing); Case Studies: General Problem Solver, Eliza, Student, Macsyma; Learning from examples; ANN and Natural Language Processing; Recent trends in AI and applications of AI algorithms. Python programming language will be used to explore and illustrate various issues and techniques in Artificial Intelligence.}
\edrow{2025}{3 \;$\bullet$\; no prerequisite and no outline published}{\begin{tight}
\item Describe core AI concepts including search algorithms, knowledge representation, and reasoning.
\item Implement basic AI algorithms for problem-solving and decision-making.
\item Analyze the applications of AI in real-world scenarios.
\item Discuss ethical considerations and limitations of AI systems.
\item Develop simple AI models using appropriate tools and frameworks.
\end{tight}}
\end{coursecard}

\begin{coursecard}{Programming for Artificial Intelligence}{the AI toolchain course}
\edrow{2017}{not specified}{\itshape The course does not exist in HEC 2017 at either level.}
\edrow{2023}{3 (2-3) \;$\bullet$\; prerequisite: Artificial Intelligence}{The first objective of the course is to introduce and then build the proficiency of students in different AI programming languages. The basics include IDE for the languages, variables, expressions, operands and operators, loops, control structures, debugging, error messages, functions, strings, lists, object-oriented constructs and basic graphics in the languages. Special emphasis is given to writing production quality clean code in the programming language. Once the classical programming languages are properly introduced, the course should introduce some libraries necessary for interpreting, analyzing and plotting numerical data in Python (e.g., NumPy, MatPlotLib, Anaconda and Pandas for Python) and give examples of each library using simple use cases and small case studies. \par\vspace{2pt}{\footnotesize\itshape Set texts: Russell, S. and Norvig, P. ``Artificial Intelligence. A Modern Approach'', 3rd ed, Prentice Hall, Inc., 2015; Norvig, P., ``Paradigms of Artificial Intelligence Programming: Case studies in Common Lisp'', Morgan Kaufman Publishers, Inc., 1992; Luger, G.F. and Stubblefield, W.A., ``AI algorithms, data structures, and idioms in Prolog, Lisp, and Java'', Pearson Addison-Wesley. 2009.}}
\edrow{2025}{3+0 \;$\bullet$\; no prerequisite and no outline published}{\begin{tight}
\item Write efficient and well-documented code.
\item Utilize version control (Git) to manage and collaborate on AI-related code repositories.
\item Implement basic Machine Learning models using standard ML frameworks
\end{tight}}
\end{coursecard}

\begin{coursecard}{Machine Learning}{the pivot of the whole cluster}
\edrow{2017}{not specified}{\itshape No undergraduate course of this name exists in HEC 2017. It is a core course of the MS Data Science programme, which publishes no outline for it.}
\edrow{2023}{3 (2-3) \;$\bullet$\; prerequisite: Artificial Intelligence}{Introduction to machine learning; concept learning: General-to-specific ordering of hypotheses, Version spaces Algorithm, Candidate elimination algorithm; Supervised Learning: decision trees, Naive Bayes, Artificial Neural Networks, Support Vector Machines, Overfitting, noisy data, and pruning, Measuring Classifier Accuracy; Linear and Logistic regression; Unsupervised Learning: Hierarchical Agglomerative Clustering. k-means partitional clustering; Self-Organizing Maps (SOM) k-Nearest-neighbor algorithm; Semi- supervised learning with EM using labeled and unlabelled data; Reinforcement Learning: Hidden Markov models, Monte Carlo inference Exploration vs. Exploitation Trade-off, Markov Decision Processes; Ensemble Learning: Using committees of multiple hypotheses. Bagging, boosting. \par\vspace{2pt}{\footnotesize\itshape Set texts: Machine Learning, Tom, M., McGraw Hill, 1997; Machine Learning: A Probabilistic Perspective, Kevin P. Murphy, MIT Press,2012.}}
\edrow{2025}{3+0 (2+1 in the SE cluster) \;$\bullet$\; no prerequisite and no outline published}{\begin{tight}
\item Explain core concepts of Machine Learning. (C2: Understand)
\item Apply machine learning algorithms to real-world datasets. (C3: Apply)
\item Analyze the performance of ML models using evaluation metrics. (C4: Analyze)
\item Implement ML solutions using modern toolkits. (C3: Apply)
\item Discuss ethical considerations while applying ML to societal problems. (A3: Valuing)
\end{tight}}
\end{coursecard}

\begin{coursecard}{Artificial Neural Networks \& Deep Learning}{the 2023 first half of deep learning}
\edrow{2017}{not specified}{\itshape The course does not exist in HEC 2017 at either level.}
\edrow{2023}{3 (2-3) \;$\bullet$\; prerequisite: Programming for Artificial Intelligence}{Introduction and history of neural networks, Basic architecture of neural networks, Perceptron and Adaline (Minimum Error Learning) for classification. Basics of deep learning, learning networks, Shallow vs. Deep learning etc.; Machine learning theory -- training and test sets, evaluation, etc. Selected topics from: Gradient descent (Delta) rule, Hebbian, Neo-Hebbian and Differential Hebbian Learning, Drive Reinforcement Theory, Kohonen Self Organizing Maps, Associative memory, Bi-directional associative memory (BAM), Energy surfaces, The Boltzmann machines, Backpropagation Networks, Feedforward Networks; Theory of Generalization; Multi-layer perceptrons, error back- propagation; Deep convolutional networks, Computational complexity of feed forward and deep convolutional neural networks; Unsupervised deep learning including auto-encoders; Deep belief networks; Restricted Boltzmann Machines; Deep Recurrent Neural Networks (BPTT, LSTM, etc.); GPU programming for deep learning CuDNN; Generative adversarial networks (GANs); Sparse coding and auto-encoders; Data augmentation, elastic distortions, data normalization; Mitigating overfitting with dropout, batch normalization, dropconnect; Novel architectures, ResNet, GoogleNet, etc \par\vspace{2pt}{\footnotesize\itshape Set texts: Neural Network Design, 2nd Edition, Martin T. Hagan, Howard, B. Demuth, Mark Hudson Beale and Orlando De Jesus, Publisher: Martin Hagan; 2 edition (September 1, 2014), ISBN-10: 0971732116; An Introduction to Neural Networks, James A Anderson, Publisher: A Bradford Book (March 16, 1995), ISBN-10: 0262011441; Fundamentals of Artificial Neural Networks, Mohammad Hassoun, Publisher: A Bradford Book (January 1, 2003), ISBN-10: 0262514672.}}
\edrow{2025}{3+0 \;$\bullet$\; no prerequisite and no outline published}{\begin{tight}
\item Explain structure, function, and types of deep neural networks.
\item Apply deep learning techniques such as CNNs, RNNs, and transfer learning for AI applications.
\item Analyze and optimize neural networks for performance using forward and backward propagation.
\item Demonstrate applications of deep reinforcement learning using modern tools.
\end{tight}}
\end{coursecard}

\begin{coursecard}{Deep Learning}{the 2023 second half, and the whole of it in 2025}
\edrow{2017}{not specified}{\itshape MS Data Science specialization core only; no outline published.}
\edrow{2023}{3 (2-3) \;$\bullet$\; prerequisite: Artificial Neural Networks and Deep Learning}{Basics of deep learning, learning networks, Shallow vs. Deep learning etc.; Machine learning theory -- training and test sets, evaluation, etc. Theory of Generalization; Multi-layer perceptrons, error back-propagation; Deep convolutional networks, Computational complexity of feed forward and deep convolutional neural networks; Unsupervised deep learning including auto-encoders; Deep belief networks; Restricted Boltzmann Machines; Deep Recurrent Neural Networks (BPTT, LSTM, etc.); GPU programming for deep learning CuDNN; Generative adversarial networks (GANs); Sparse coding and auto-encoders; Data augmentation, elastic distortions, data normalization; Mitigating overfitting with dropout, batch normalization, dropconnect; Novel architectures, ResNet, GoogleNet, etc \par\vspace{2pt}{\footnotesize\itshape Set texts: Deep Learning by Ian Goodfellow, Yoshua Bengio, Aaron Courville (; Deep learning with python by Francoise Chollet, ISBN-10: 9781617294433, 2017.}}
\edrow{2025}{3+0 \;$\bullet$\; no prerequisite and no outline published}{\begin{tight}
\item Understand the fundamentals of neural networks in AI C2 (Understand)
\item Explain how simple ANNs can be designed C2 (Understand)
\item Apply deep learning algorithms to real-world classification problems C3 (Apply)
\item Analyze results from deep learning to select appropriate solutions C4 (Analyze)
\end{tight}}
\end{coursecard}

\begin{coursecard}{Introduction to Data Science}{the process-and-EDA course}
\edrow{2017}{not specified}{\itshape The nearest equivalent, Tools and Techniques in Data Science, is an MS core course with no published outline.}
\edrow{2023}{3 (2-3) \;$\bullet$\; prerequisite: Artificial Intelligence}{Introduction: What is Data Science? Big Data and Data Science hype, Datafication, Current landscape of perspectives, Skill sets needed; Statistical Inference: Populations and samples, Statistical modeling, probability distributions, fitting a model, Intro to Python; Exploratory Data Analysis and the Data Science Process; Basic Machine Learning Algorithms: Linear Regression, k-Nearest Neighbors (k-NN), k-means, Naive Bayes; Feature Generation and Feature Selection; Dimensionality Reduction: Singular Value Decomposition, Principal Component Analysis; Mining Social-Network Graphs: Social networks as graphs, Clustering of graphs, Direct discovery of communities in graphs, Partitioning of graphs, Neighborhood properties in graphs; Data Visualization: Basic principles, ideas and tools for data visualization; Data Science and Ethical Issues: Discussions on privacy, security, ethics, Next-generation data scientists. \par\vspace{2pt}{\footnotesize\itshape Set texts: Foundations of data science, Blum, A., Hopcroft, J., \& Kannan, R., Vorabversion eines Lehrbuchs, 2016; An Introduction to Data Science, Jeffrey S. Saltz, Jeffrey M. Stanton, SAGE Publications, 2017; Python for everybody: Exploring data using Python 3, Severance, C.R., CreateSpace Independent Pub Platform. 2016.}}
\edrow{2025}{3+0 \;$\bullet$\; no prerequisite and no outline published}{\begin{tight}
\item Understand the fundamental principles of data science. C2 (Understand)
\item Apply EDA and the Data Science processes. C3 (Apply)
\item Apply basic machine learning algorithms to solve real-world problems of moderate complexity. C3 (Apply)
\end{tight}}
\end{coursecard}

\begin{coursecard}{Data Mining}{pattern discovery}
\edrow{2017}{not specified}{\itshape The course does not exist at any level in HEC 2017.}
\edrow{2023}{3 (2-3) \;$\bullet$\; prerequisite: Advance Statistics, Introduction to Data Science}{Introduction to data mining and basic concepts, Pre-Processing Techniques \& Summary Statistics, Association Rule mining using Apriori Algorithm and Frequent Pattern Trees, Introduction to Classification Types, Supervised Classification (Decision trees, Naïve Bayes Classification, K-Nearest Neighbors, Support Vector Machines etc.), Unsupervised Classification (K Means, K Median, Hierarchical and Divisive Clustering, Kohonen Self Organizing maps), outlier \& anomaly detection, Web and Social Network Mining, Data Mining Trends and Research Frontiers. Implementing concepts using Python \par\vspace{2pt}{\footnotesize\itshape Set texts: Jiawei Han \& Micheline Kamber, Jian Pei (2011). Data Mining: Concepts and Techniques, 3rd Edition; Pang-Ning Tan, Michael Steinbach, and Vipin Kumar (2005). Introduction to Data Mining; Charu C. Aggarwal (2015). Data Mining: The Textbook.}}
\edrow{2025}{3+0 \;$\bullet$\; no prerequisite and no outline published}{\begin{tight}
\item Explain fundamentals of data mining concepts, techniques, and algorithms.C2 (Understand)
\item Apply data mining algorithms to real-world datasets using appropriate tools and techniques.C3 (Apply)
\item Evaluate and interpret the performance of data mining models for informed decision-making.C5 (Evaluate) 3. ARTIFICIAL INTELLIGENCE ELECTIVES:
\end{tight}}
\end{coursecard}

\begin{coursecard}{Big Data Analytics}{the distributed compute course}
\edrow{2017}{3(2,1) \;$\bullet$\; prerequisite: Probability and Statistics, Programming Fundamentals}{Introduction to Big Data Analytics, Big Data Platforms, Data Store \& Processing using Hadoop, Big Data Storage and Analytics, Big Data Analytics ML Algorithms, Recommendation, Clustering, and Classification, Linked Big Data: Graph Computing and Graph Analytics, Graphical Models and Bayesian Networks, Big Data Visualization, Cognitive Mobile Analytics. \par\vspace{2pt}{\footnotesize\itshape Set texts: Mining of Massive Datasets, Jure Leskovec, Anand Rajaraman, Jeff Ullman, 2nd edition, 2011; Hadoop: The Definitive Guide, Tom White, 4th edition. 2009; Data-Intensive Text Processing with Map Reduce, Jimmy Lin and Chris, 2010.}}
\edrow{2023}{3 (2-3) \;$\bullet$\; prerequisite: Introduction to Data Science}{Introduction and Overview of Big Data Systems; Platforms for Big Data, Hadoop as a Platform, Hadoop Distributed File Systems (HDFS), MapReduce Framework, Resource Management in the cluster (YARN), Apache Scala Basic, Apache Scala Advances, Resilient Distributed Datasets (RDD), Apache Spark, Apache Spark SQL, Data analytics on Hadoop / Spark, Machine learning on Hadoop / Spark, Spark Streaming, Other Components of Hadoop Ecosystem \par\vspace{2pt}{\footnotesize\itshape Set texts: White, Tom. ``Hadoop: The definitive guide." O'Reilly Media, Inc., 2012; Karau, Holden, Andy Konwinski, Patrick Wendell, and Matei Zaharia. ``Learning spark: lightning-fast big data analysis." O'Reilly Media, Inc., 2015; Miner, Donald, and Adam Shook. ``MapReduce design patterns: building effective algorithms and analytics for Hadoop and other systems." O'Reilly Media, Inc., 2012..}}
\edrow{2025}{3+0 \;$\bullet$\; no prerequisite and no outline published}{\begin{tight}
\item Understand the fundamental concepts of Big Data and its programming paradigm. C3(Understand)
\item Hadoop/MapReduce Programming, Framework, and Ecosystem C2 (Understand)
\item Apache Spark Programming C3 (Apply)
\end{tight}}
\end{coursecard}

\begin{coursecard}{Advanced / Computational Statistics}{the statistical support course}
\edrow{2017}{not specified}{\itshape The course does not exist in HEC 2017 at either level.}
\edrow{2023}{3 (3-0) \;$\bullet$\; prerequisite: Probability and Statistics}{Introduction to Statistics, Use of Statistics in Data Science, Experimental Design, Statistical Techniques for Forecasting, Interpolation/ Extrapolation, Introduction to Probability, Conditional Probability, Prior and Posterior Probability, Random number generation (RNG), Techniques for RNG, Correlation analysis, Chi Square Dependency tests, Diversity Index, Data Distributions Multivariate Distributions, Error estimation, Confidence Intervals, Linear transformations, Gradient Descent and Coordinate Descent, Likelihood inference, Revision of linear regression and likelihood inference, Fitting algorithms for nonlinear models and related diagnostics, Generalized linear model; exponential families; variance and link functions, Proportion and binary responses; logistic regression, Count data and Poisson responses; log-linear models, Overdispersion and quasi-likelihood; estimating functions, Mixed models, random effects, generalized additive models and penalized regression; Introduction to SPSS, Probability/ Correlation analysis/ Dependency tests/ Regression in SPSS. \par\vspace{2pt}{\footnotesize\itshape Set texts: Probability and Statistics for Computer Scientists, 2nd Edition, Michael Baron; Probability for Computer Scientists, online Edition, David Forsyth; Discovering Statistics using SPSS for Windows, Andy Field.}}
\edrow{2025}{3+0 \;$\bullet$\; no prerequisite and no outline published}{\begin{tight}
\item Understand the fundamental statistics for data science and applications of statistics in data science.C1 (Knowledge)
\item Apply Statistical techniques in real-life problems. C3(Apply)
\item Apply basic data science statistical techniques by using tools on real-world datasets. C3 (Apply)
\end{tight}}
\end{coursecard}

\begin{coursecard}{Data Warehousing \& Business Intelligence}{the enterprise data stack}
\edrow{2017}{not specified}{\itshape The course does not exist in HEC 2017 at either level.}
\edrow{2023}{2-1 \;$\bullet$\; prerequisite: Introduction to Data Science}{Introduction to Data Warehouse and Business Intelligence; Necessities and essentials of Business Intelligence; DW Life Cycle and Basic Architecture; DW Architecture in SQL Server; Logical Model; Indexes; Physical Model; Optimizations; OLAP Operations, Queries and Query Optimization; Building the DW; Data visualization and reporting based on Datawarehouse using SSAS and Tableau; Data visualization and reporting based on Cube; Reports and Dashboard management on PowerBI; Dashboard Enrichment; Business Intelligence Tools. \par\vspace{2pt}{\footnotesize\itshape Set texts: W. H. Inmon, ``Building the Data Warehouse'', Wiley-India Edition; Ralph Kimball, ``The Data Warehouse Toolkit -- Practical Techniques for Building Dimensional Data Warehouse,'' John Wiley \& Sons, Inc; Matteo Golfarelli, Stefano Rizzi, ``Data Warehouse Design - Modern Principles and Methodologies'', McGraw Hill Publisher.}}
\edrow{2025}{3+0 \;$\bullet$\; no prerequisite and no outline published}{\begin{tight}
\item Explain the concepts, architecture, and components of Business Intelligence systems.C2 (Understand)
\item Apply BI tools and techniques to extract, transform, and visualize data for decision- making.C3 (Apply)
\item Design and evaluate BI solutions to support strategic and operational business decisions.. C3 (Evaluate)
\end{tight}}
\end{coursecard}

\begin{coursecard}{Data Visualization}{exploratory graphics}
\edrow{2017}{not specified}{\itshape The course does not exist in HEC 2017 at either level.}
\edrow{2023}{3 (2-3) \;$\bullet$\; prerequisite: Data Warehousing \& Business Intelligence}{Introduction of Exploratory Data Analysis and Visualization, Building Blocks and Basic Operations; Types of Exploratory Graphs, single and multi-dimensional summaries, five number summary, box plots, histogram, bar plot and others; Distributions, their representation using histograms, outliers, variance; Probability Mass Functions and their visualization; Cumulative distribution functions, percentile-based statistics, random numbers; Modelling distributions, exponential, normal, lognormal, pareto; Probability density functions, kernel density estimation; Relationship between variables, scatter plots, correlation, covariance; Estimation and Hypothesis Testing; Clustering using K-means and Hierarchical; Time series and survival analysis; Implementing concepts with R (or similar language) \par\vspace{2pt}{\footnotesize\itshape Set texts: ``Exploratory Data Analysis with R'' by Roger D. Peng.}}
\edrow{2025}{3+0 \;$\bullet$\; no prerequisite and no outline published}{\begin{tight}
\item Describes the basic concept and techniques for any given raw data visualization. C3 (Understand)
\item Select and apply a proper data mining algorithm to discover interesting patterns C3 (Apply)
\item Analyze and extract patterns to solve problems and point out how to deploy solution C4 (Analyze)
\item Evaluate systematically supervised, semi-supervised, and unsupervised models and algorithms with respect to their accuracy C4 (Analyze)
\end{tight}}
\end{coursecard}

\begin{coursecard}{Natural Language Processing}{the same title, three different courses}
\edrow{2017}{3(3,0) \;$\bullet$\; prerequisite: none}{Deterministic and stochastic grammars, Parsing algorithms, CFGs, Representing meaning / Semantics, Semantic roles, Temporal representations, Corpus-based methods, N-grams and HMMs, Smoothing and backoff, POS tagging and morphology, Information retrieval, Vector space model, Precision and recall, Information extraction, Language translation, Text classification, categorization, Bag of words model. \par\vspace{2pt}{\footnotesize\itshape Set texts: Python Machine Learning, Sebastian Raschka. Publisher: Packt Publishing, 2015; Natural Language Processing with Python: Analyzing Text with the Natural Language Toolkit Latest Edition, Steven Bird, Ewan Klein and Edward Loper Publisher: O'Reilly Media, 2009; Speech and Language Processing, Latest Edition, Daniel Jurafsky and James H. Martin Publisher: Prentice Hall, 2000.}}
\edrow{2023}{3 (2-3) \;$\bullet$\; prerequisite: Artificial Neural Networks \& Deep Learning}{Introduction \& History of NLP, Parsing algorithms, Basic Text Processing, Minimum Edit Distance, Language Modeling, Spelling Correction, Text Classification, Deterministic and stochastic grammars, CFGs, Representing meaning /Semantics, Semantic roles, Semantics and Vector models, Sentiment Analysis, Temporal representations, Corpus-based methods, N-grams and HMMs, Smoothing and backoff, POS tagging and morphology, Information retrieval, Vector space model, Precision and recall, Information extraction, Relation Extraction (dependency, constituency grammar), Language translation, Text classification, categorization, Bag of words model, Question and Answering, Text Summarization \par\vspace{2pt}{\footnotesize\itshape Set texts: Daniel Jurafsky and James H. Martin. 2018. Speech and Language Processing: An Introduction to Natural Language Processing. Third Edition. Prentice Hall; Foundations of Statistical Natural Language Processing, Manning and Schütze, MIT Press. Cambridge, MA: May 1999.}}
\edrow{2025}{3+0 \;$\bullet$\; no prerequisite and no outline published}{\begin{tight}
\item Explain the fundamental concepts and techniques of Natural Language Processing, including tokenization, POS tagging, and parsing. C2 (Understand)
\item Apply NLP libraries and frameworks (e.g., NLTK ) to preprocess and analyze text data for real-world applications.C3 (Apply)
\item Design and implement NLP-based solutions such as sentiment analysis or text classification using machine learning models. C6 (Create)
\end{tight}}
\end{coursecard}

\begin{coursecard}{Computer Vision}{images and recognition}
\edrow{2017}{3 \;$\bullet$\; prerequisite: none}{Introduction, Image formation, Spatial and frequency domain processing, Feature detection and extraction, Image registration, Segmentation, Camera calibration, Structure from motion, Motion estimation, Stereo vision, Object detection and recognition, Object tracking, 3D scene reconstruction, Context and scene understanding, Image stitching, Image-based and video-based rendering, High-performance computing paradigms for vision and image processing., \par\vspace{2pt}{\footnotesize\itshape Set texts: Computer Vision - A Modern Approach, by D. Forsyth and J. Ponce, Prentice Hall, 2003; Szeliski R., Computer Vision - Algorithms and Applications, Springer, 2011; J. R. Parker, Algorithms for Image Processing and Computer Vision, Willey Publishing Inc. 2011.}}
\edrow{2023}{3 (2-3) \;$\bullet$\; prerequisite: Artificial Neural Networks \& Deep Learning}{Introduction to Computer Vision (Problems faced, History and Modern Advancements). Image Processing, Image filtering, Image pyramids and Fourier transform, Hough transform. Camera models, Setting up a camera model from parameters, Camera looking at a plane, Relationship of plane and horizon line, Rotation about camera center. Concatenation, Decomposition and Estimation of transformation from point correspondences, Points and planes in 2D/3D, Transformations in 2D/3D, Rotations in 2D/3D. Edge detection, corner detection. Feature descriptors and matching (HoG features, SIFT, SURF). Applications of Computer Vision Traditional Methods: Image Stitching: Making a bigger picture from smaller pictures Single View Geometry: Converting a single image into a 3D model. Applications of CV using Deep Learning: Image Detection (Localization, Historical Techniques, RCNN, FRCNN, YOLO, Retina), Image Segmentation (UNet, SegNet, MaskRCNN), Image Generation (GANN) \par\vspace{2pt}{\footnotesize\itshape Set texts: Computer Vision: Algorithms and Applications, by Richard Szeliski. Reference Book:; Multiple View Geometry in Computer Vision, by Richard Hartley and Andrew Zisserman; Computer Vision: A Modern Approach, by David Forsyth and Jean Ponce.}}
\edrow{2025}{3+0 \;$\bullet$\; no prerequisite and no outline published}{\begin{tight}
\item Understand and explain the fundamental concepts, techniques, and algorithms used in Computer Vision, including image processing, feature extraction, and object recognition.C2 (Understand)
\item Implement computer vision algorithms and techniques using modern libraries and frameworks, such as OpenCV, TensorFlow, or PyTorch, to solve real-world image processing problems.C3 (Apply)
\item Evaluate and assess the performance of computer vision models using metrics like accuracy, precision, recall, Intersection over Union (IoU), and mean average precision (mAP).C4 (Analyze)
\end{tight}}
\end{coursecard}

\begin{coursecard}{Knowledge Representation \& Reasoning}{the symbolic half of AI}
\edrow{2017}{not specified}{\itshape The course does not exist in HEC 2017 at either level.}
\edrow{2023}{3 (3-0) \;$\bullet$\; prerequisite: Artificial Intelligence}{Propositional Logic, First-order Logic, Horn Clauses, Description Logic, Reasoning using Description Logic, Forward and Backward Chaining in Inference Engines, Semantic Networks, Ontologies and Ontology Languages, Logical Agents, Planning, Rule-based Knowledge Representation, Reasoning Under Uncertainty, Bayesian Networks Representation, Inference in Bayesian Networks, Fuzzy Logic, Inferene using Fuzzy Rules, Markov Models, Commonsense Reasoning, Explainable AI. \par\vspace{2pt}{\footnotesize\itshape Set texts: Stuard Russell and Peter Norvig, Artificial Intelligence: A Modern Approach (3rd Ed.) (2015); David Poole and Alan Mackworth, Artificial Intelligence: Foundations of Computational Agents, 2nd Ed, 2017; Ronald Brachman and Hector Levesque. Knowledge Representation and Reasoning, 2004.}}
\edrow{2025}{3+0 \;$\bullet$\; no prerequisite and no outline published}{\begin{tight}
\item Explain the fundamental concepts, models and techniques of knowledge representation and reasoning.
\item Apply logical formalism, ontologies and inference mechanisms
\item Use appropriate tools and methods to implement reasoning systems and evaluate their effectiveness.
\end{tight}}
\end{coursecard}

\begin{coursecard}{Reinforcement Learning}{listed in 2023, specified only in 2025}
\edrow{2017}{not specified}{\itshape The course does not exist in HEC 2017 at either level.}
\edrow{2023}{not specified}{\itshape Listed as a BSAI domain elective but not among the 46 courses for which HEC 2023 publishes contents.}
\edrow{2025}{3+0 \;$\bullet$\; no prerequisite and no outline published}{\begin{tight}
\item Understand the core components of Markov decision processes and Reinforcement Learning
\item Implement reinforcement learning algorithms to model and solve dynamic environments
\item Evaluate reinforcement learning approaches in terms of performance scalability and applicability to real world domains
\end{tight}}
\end{coursecard}

\begin{coursecard}{Speech Processing}{listed in 2023, specified only in 2025}
\edrow{2017}{not specified}{\itshape The course does not exist in HEC 2017 at either level.}
\edrow{2023}{not specified}{\itshape Listed as a BSAI and BSDS domain elective but not among the 46 specified courses.}
\edrow{2025}{3+0 \;$\bullet$\; no prerequisite and no outline published}{\begin{tight}
\item Explain fundamental concepts of speech production, perception, and signal processing, and apply them to analyze and model speech signals.
\item Implement and evaluate speech processing techniques using modern programming tools and frameworks.
\item Evaluate effectiveness of speech processing techniques along with ethical and societal implications of speech technologies such as privacy, surveillance, and accessibility.
\end{tight}}
\end{coursecard}

\begin{coursecard}{Tools and Techniques in Data Science}{an MS core course in 2017, an undergraduate elective in 2025}
\edrow{2017}{not specified}{\itshape The course does not exist in HEC 2017 at either level.}
\edrow{2023}{not specified}{\itshape Not among the 46 courses for which HEC 2023 publishes contents.}
\edrow{2025}{3+0 \;$\bullet$\; no prerequisite and no outline published}{\begin{tight}
\item Explain the purpose and functionality of essential tools and techniques used in data science workflows.C2 (Understand)
\item Apply data preprocessing, transformation, and visualization techniques using appropriate tools to solve practical problems..C3 (Apply)
\item Design and implement integrated data science workflows using multiple tools for real-world datasets.C6(Create)
\end{tight}}
\end{coursecard}

\begin{coursecard}{Data Engineering}{new in 2025}
\edrow{2017}{not specified}{\itshape The course does not exist in HEC 2017 at either level.}
\edrow{2023}{not specified}{\itshape Not among the 46 courses for which HEC 2023 publishes contents.}
\edrow{2025}{3+0 \;$\bullet$\; no prerequisite and no outline published}{\begin{tight}
\item Explain the principles of data engineering, including data pipelines, ETL processes, and storage systems. C3 (Understand)
\item Apply data integration and transformation techniques to build scalable and efficient data pipelines. C3 (Apply)
\item Design and implement data architectures for big data and real-time analytics using modern tools and platforms. C3 (Apply)
\end{tight}}
\end{coursecard}

\begin{coursecard}{Generative AI}{new in 2025}
\edrow{2017}{not specified}{\itshape The course does not exist in HEC 2017 at either level.}
\edrow{2023}{not specified}{\itshape Not among the 46 courses for which HEC 2023 publishes contents.}
\edrow{2025}{3+0 \;$\bullet$\; no prerequisite and no outline published}{\begin{tight}
\item Describe the foundational concepts, models, and architectures of Generative AI, including GANs, VAEs, and large language models.C2 (Understand)
\item Implement and fine-tune generative models using modern AI frameworks such as TensorFlow or PyTorch. C3 (Apply)
\item Analyze the ethical, societal, and practical implications of deploying Generative AI in real-world applications. C4 (Analyze)
\end{tight}}
\end{coursecard}

\begin{coursecard}{Agentic AI}{new in 2025}
\edrow{2017}{not specified}{\itshape The course does not exist in HEC 2017 at either level.}
\edrow{2023}{not specified}{\itshape Not among the 46 courses for which HEC 2023 publishes contents.}
\edrow{2025}{3+0 \;$\bullet$\; no prerequisite and no outline published}{\begin{tight}
\item Demonstrate understanding of agentic AI concepts and design autonomous AI agents capable of reasoning, planning, and tool-use for real-world tasks.
\item Implement agentic AI systems using modern frameworks with integration of APIs, databases, and external tools.
\item Evaluate agentic AI systems with special focus on ethical, societal, and professional implications of deploying autonomous AI agents, and propose responsible design strategies.
\end{tight}}
\end{coursecard}

\begin{coursecard}{Machine Learning Operations (MLOps)}{new in 2025, and where the two lanes meet}
\edrow{2017}{not specified}{\itshape The course does not exist in HEC 2017 at either level.}
\edrow{2023}{not specified}{\itshape Not among the 46 courses for which HEC 2023 publishes contents.}
\edrow{2025}{3+0 \;$\bullet$\; no prerequisite and no outline published}{\begin{tight}
\item Explain concepts of MLOps and model lifecycle management.
\item Apply automation for data pipelines, training, and deployment.
\item Evaluate reproducibility, scalability, and governance in ML systems.
\end{tight}}
\end{coursecard}

\begin{coursecard}{Information Retrieval}{new as a course in 2025}
\edrow{2017}{not specified}{\itshape The course does not exist in HEC 2017 at either level.}
\edrow{2023}{not specified}{\itshape Not among the 46 courses for which HEC 2023 publishes contents.}
\edrow{2025}{3+0 \;$\bullet$\; no prerequisite and no outline published}{\begin{tight}
\item Explain the core principles and algorithms used in Information Retrieval systems, including indexing, ranking, and query processing.C2 (Understand)
\item Apply IR techniques to build a simple search engine or retrieval system, using tools and programming languages like Python, Elasticsearch, or Apache Lucene.C3 (Apply)
\item Evaluate the performance of Information Retrieval systems through metrics such as precision, recall, F1-score, and Mean Average Precision (MAP) C4 (Analyze)
\end{tight}}
\end{coursecard}

\begin{coursecard}{Data Ethics \& Security}{new in 2025}
\edrow{2017}{not specified}{\itshape The course does not exist in HEC 2017 at either level.}
\edrow{2023}{not specified}{\itshape Not among the 46 courses for which HEC 2023 publishes contents.}
\edrow{2025}{3+0 \;$\bullet$\; no prerequisite and no outline published}{\begin{tight}
\item Explain the principles of data ethics, privacy laws, and governance frameworks.C2 (Understand)
\item Analyze ethical and legal implications of data collection, storage, and sharing in real- world scenarios.C4 (Analyze)
\item Evaluate compliance strategies for data protection and governance in organizational contexts.C5 (Evaluate)
\item Design policies and guidelines to ensure ethical data usage and regulatory compliance.C6 (Create)
\end{tight}}
\end{coursecard}

\begin{coursecard}{Optimization Techniques}{new in 2025}
\edrow{2017}{not specified}{\itshape The course does not exist in HEC 2017 at either level.}
\edrow{2023}{not specified}{\itshape Not among the 46 courses for which HEC 2023 publishes contents.}
\edrow{2025}{3+0 \;$\bullet$\; no prerequisite and no outline published}{\begin{tight}
\item Explain the theoretical basis of optimization problems.
\item Formulate and solve optimization problems using analytical and computational methods.
\item Implement and evaluate most effective optimization tools and techniques for real world problems.
\end{tight}}
\end{coursecard}

\begin{coursecard}{Stochastic Processes}{in 2017 and 2025, absent from 2023}
\edrow{2017}{3 (3,0) \;$\bullet$\; prerequisite: Probability and Statistics}{Discrete Markov chains, classification of states, first passage and recurrence times, absorption problems, stationary and limiting distributions. Chapman-Kolmogorov equations, Long run behavior of Markov chains, Absorption probabilities and expected times to absorption, Statistical aspects of Markov chains, The mover-stayer model, Application of a Markov chain and mover-stayer model to modeling repayment behavior of bank loans' grantees. Markov Processes in continuous time: Poisson processes, birth- death processes. Poisson process The Kolmogorov differential equations, Limiting behavior of continuous time Markov chains The Q matrix, forward and backward differential equations, imbedded Markov Chain, stationary distribution. renewal theory, Brownian Motion and its generalizations, Discrete time martingales, Conditional expectation, Definition of a martingale and examples, Optional stopping theorem, Stochastic calculus \par\vspace{2pt}{\footnotesize\itshape Set texts: Introduction to Probability Models, 11th Ed, Sheldon M. Ross, Academic Press 2014; Essentials of stochastic processes, Durrett, Richard. Springer Science \& Business Media, 2nd Ed, 2012; Introduction to Stochastic Processes, 2nd Ed, G.F. Lawler, Chapman and Hall, Probability Series, 2006.}}
\edrow{2023}{not specified}{\itshape The course is absent from HEC 2023 altogether.}
\edrow{2025}{3+0 \;$\bullet$\; no prerequisite and no outline published}{\begin{tight}
\item Explain probability theory and stochastic models to characterize randomness and uncertainty in computational systems
\item Apply probabilistic techniques and stochastic methods to analyze and model real world data and processes
\item Evaluate the suitability and limitations of stochastic models in engineering and computer science applications
\end{tight}}
\end{coursecard}


\subsection{The infrastructure courses}

\begin{coursecard}{Computer Networks}{compulsory in every degree of every edition}
\edrow{2017}{3+1 \;$\bullet$\; prerequisite: none}{Introduction and protocols architecture, basic concepts of networking, network topologies, layered architecture, physical layer functionality, data link layer functionality, multiple access techniques, circuit switching and packet switching, LAN technologies, wireless networks, MAC addressing, networking devices, network layer protocols, IPv4 and IPv6, IP addressing, subnetting, CIDR, routing protocols, transport layer protocols, ports and sockets, connection establishment, flow and congestion control, application layer protocols, latest trends in computer networks. \par\vspace{2pt}{\footnotesize\itshape Set texts: Computer Networking: A Top-Down Approach Featuring the Internet, 6th edition by James F. Kurose and Keith W. Ross; Computer Networks, 5th Edition by Andrew S. Tanenbaum; Data and Computer Communications, 10th Edition by William Stallings.}}
\edrow{2023}{3 (2-3) \;$\bullet$\; prerequisite: none}{Introduction and protocols architecture, basic concepts of networking, network topologies, layered architecture, physical layer functionality, data link layer functionality, multiple access techniques, circuit switching and packet switching, LAN technologies, wireless networks, MAC addressing, networking devices, network layer protocols, IPv4 and IPv6, IP addressing, subnetting, CIDR, routing protocols, transport layer protocols, ports and sockets, connection establishment, flow and congestion control, application layer protocols, latest trends in computer networks.}
\edrow{2025}{3 \;$\bullet$\; no prerequisite and no outline published}{\begin{tight}
\item Explain the fundamental concepts of computer networking, including protocols, topologies, and models (OSI, TCP/IP).
\item Configure and troubleshoot basic network devices and connections.
\item Analyze network security threats and mitigation techniques.
\item Demonstrate understanding of data transmission and error handling.
\item Design simple network architectures to meet organizational needs.
\end{tight}}
\end{coursecard}

\begin{coursecard}{System \& Network Administration}{compulsory, compulsory, then optional}
\edrow{2017}{3 (3,0) \;$\bullet$\; prerequisite: Operating System}{Introduction To System Administration. SA Components. Server Environment (Microsoft and Linux). Reliable Products, Server Hardware Costing, Maintenance Contracts and Spare Parts, Maintaining Data Integrity, Client Server OS Configuration, Providing Remote Console Access. Comparative Analysis of OS: Important Attributes, Key Features, Pros and Cons. Linux Installation and Verification, Configuring Local Services and Managing Basic System Issues. Administer Users and Groups. Software Management. Managing Network Services and Network Monitoring Tools. Boot Management and Process Management. IP Tables and Filtering. Securing Network Traffic. Advanced File Systems and Logs. Bash Shell Scripting. Configuring Servers (FTP, NFS, Samba, DHCP, DNS and Apache). \par\vspace{2pt}{\footnotesize\itshape Set texts: The Practice of System and Network Administration, Second Edition by Thomas Limoncelli, Christina Hogan and Strata Chalup, Addison-Wesley Professional; 2nd Edition (2007). ISBN-10: 0321492668; Red Hat Enterprise Linux 6 Bible: Administering Enterprise Linux Systems by William vonHagen, 2011; Studyguide for Practice of System and Network Administration by Thomas A. Limoncelli, Cram101; 2nd Edition (2011). ISBN-10: 1428851755.}}
\edrow{2023}{not specified}{\itshape A BSIT domain core course in Semester 5, and unspecified for the same reason: HEC 2023 has no Information Technology course-contents table.}
\edrow{2025}{2+1 \;$\bullet$\; no prerequisite and no outline published}{\begin{tight}
\item Explain principles of system and network administration.
\item Configure and troubleshoot network services and operating systems.
\item Use administrative tools to monitor and secure systems.
\item Demonstrate professionalism in managing shared computing resources.
\item Work effectively in teams to maintain IT infrastructure.
\end{tight}}
\end{coursecard}

\begin{coursecard}{Cloud Computing}{unspecified, unspecified, then compulsory}
\edrow{2017}{not specified}{\itshape Named once, in the BS Software Engineering elective list, at 3-0. It is absent from the 75-course bachelor list and no specification exists anywhere in the booklet.}
\edrow{2023}{not specified}{\itshape Listed as a BSIT domain elective and in the master course list, but HEC 2023 publishes contents only for courses common to all computing degrees and for the AI, Data Science and Cyber Security clusters. There is no Information Technology table.}
\edrow{2025}{3 \;$\bullet$\; no prerequisite and no outline published}{\begin{tight}
\item Describe fundamental cloud computing models and services (IaaS, PaaS, SaaS).
\item Deploy and manage applications in cloud environments.
\item Analyze the benefits and challenges of cloud computing.
\item Implement basic cloud security and compliance measures.
\item Evaluate cloud solutions for scalability, cost, and performance.
\end{tight}}
\end{coursecard}

\begin{coursecard}{Cloud Infrastructure and Services}{the 2025 IT-cluster near-duplicate of Cloud Computing}
\edrow{2017}{not specified}{\itshape The course does not exist in HEC 2017 at either level.}
\edrow{2023}{not specified}{\itshape Not among the 46 courses for which HEC 2023 publishes contents.}
\edrow{2025}{2+1 \;$\bullet$\; no prerequisite and no outline published}{\begin{tight}
\item Explain cloud service models and infrastructure components.
\item Deploy and manage virtualized cloud environments.
\item Use cloud management tools for provisioning and monitoring.
\item Evaluate legal and compliance issues in cloud computing.
\item Engage in lifelong learning to adapt to cloud innovations.
\end{tight}}
\end{coursecard}

\begin{coursecard}{Virtual Systems and Services}{2017's nearest thing to a cloud course}
\edrow{2017}{3 \;$\bullet$\; prerequisite: Programming Fundamentals}{This course will investigate the current state of virtualization in computing systems. Virtualization at both the hardware and software levels will be examined, with emphasis on the hypervisor configurations of systems such as Vmware, Zen and Hyper-V. The features and limitations of virtual environments will be considered, along with several case studies used to demonstrate the configuration and management of such systems. Para-virtualized software components will be analyzed and their pros and cons discussed. Processor and peripheral support for virtualization will also be examined, with a focus on emerging hardware features and the future of virtualization. \par\vspace{2pt}{\footnotesize\itshape Set texts: Handbook of Virtual Environments: Design, Implementation, and Applications (Human Factors and Ergonomics), Edited by Kay M Stanney, Lawrence Erlbaum Associates Virtual Reality Technology by GRIGORE.}}
\edrow{2023}{not specified}{\itshape A BSIT domain elective example, unspecified: HEC 2023 has no Information Technology course-contents table.}
\edrow{2025}{not specified}{\itshape The course is gone. Its ground is covered by Cloud Computing and, in the infrastructure cluster, by Microservices Architecture and Docker Containers.}
\end{coursecard}

\begin{coursecard}{Information Technology Infrastructure}{the architecture-side companion to administration}
\edrow{2017}{3 (3,0) \;$\bullet$\; prerequisite: none}{Definition of IT Infrastructure, Non-functional Attributes, Availability Concepts, Sources of Unavailability, Availability Patterns. Performance. Security Concepts. Data centres. Servers: Availability, Performance, Security. Networking: Building Blocks, Availability, Performance, Security. Storage: Availability, Performance, Security. Virtualization: Availability, Performance, Security. Operating Systems: Building Blocks, Implementing Various OSs, OS availability, OS Performance, OS Security. End User Devises: Building Blocks, Device Availability, Performance, Security. IT Infrastructure Management. Service Delivery Processes. Service Support Processes. Ethics, Trends, organizational and technical issues related to IT infrastructure. \par\vspace{2pt}{\footnotesize\itshape Set texts: IT Infrastructure Architecture: Infrastructure building blocks and concepts by Sjaak Laan, Lulu.com (November 5, 2011). ISBN-10: 1447881281; IT Infrastructure and its Management by Prof Phalguni Gupta, Tata McGraw Hill Education Private Limited (October 6, 2009). ISBN-10: 0070699798; IT Architecture for Dummies by Kalani Kirk Hausman and Susan Cook, For Dummies; 1st Edition (November 9, 2010). ISBN-10: 0470554231.}}
\edrow{2023}{not specified}{\itshape A BSIT domain core course, unspecified for the same reason.}
\edrow{2025}{2+1 \;$\bullet$\; no prerequisite and no outline published}{\begin{tight}
\item Identify components of IT infrastructure and their roles.
\item Design basic infrastructure models for organizational needs.
\item Apply tools for infrastructure monitoring and support.
\item Assess societal and organizational impacts of infrastructure decisions.
\item Engage in self-directed learning to stay updated with infrastructure trends.
\end{tight}}
\end{coursecard}

\begin{coursecard}{Network Security}{specified in 2023 only because it is common to every degree}
\edrow{2017}{not specified}{\itshape The course does not exist in HEC 2017 at either level.}
\edrow{2023}{3 (2-3) \;$\bullet$\; prerequisite: Introduction to Cyber Security}{Introduction to network security, Networking Concepts and Protocols, Network Threats and Vulnerabilities, Network Security Planning and Policy, Access Control, Defense against Network Attacks, DOS and DDOS detection and prevention, Firewalls, Intrusion Detection and Prevention Systems, Antivirus Filtering, Naming and DNS Security, DNSSEC, IP security, Secure Sockets Layer, VPN, Packet Sniffing and spoofing, Honeypot, Ethernet Security, Wireless Security, Wireless Attacks, Wireless LAN Security with 802.11i, Wireless Security Protocols, Wireless Intrusion Detection, Physical access and Security, Tor Network, Network Forensics. Defense against Network Attacks. \par\vspace{2pt}{\footnotesize\itshape Set texts: Network Security Assessment: Know Your Network by Chris McNab, 3rd Edition or latest; Corporate Computer Security, by Randall J. Boyle, 3th Edition; Bulletproof Wireless Security by Praphul Chandra.}}
\edrow{2025}{2+1 \;$\bullet$\; no prerequisite and no outline published}{\begin{tight}
\item Explain common network threats and vulnerabilities. (BT2 -- Understand) (PLO1)
\item Configure basic network security mechanisms such as firewalls and VPNs. (BT3 -- Apply) (PLO5)
\item Design and develop network defense solutions (BT3-Apply)(PLO3)
\item Evaluate existing security architectures to identify strengths and weaknesses. (BT5 -- Evaluate) (PLO1)
\end{tight}}
\end{coursecard}

\begin{coursecard}{Parallel and Distributed Computing}{a domain core in BSAI, BSDS and BSIT alike}
\edrow{2017}{not specified}{\itshape The course does not exist in HEC 2017 at either level.}
\edrow{2023}{3 (2-3) \;$\bullet$\; prerequisite: Object Oriented Programming, Operating Systems}{Asynchronous/synchronous computation/communication, concurrency control, fault tolerance, GPU architecture and programming, heterogeneity, interconnection topologies, load balancing, memory consistency model, memory hierarchies, Message passing interface (MPI), MIMD/SIMD, multithreaded programming, parallel algorithms \& architectures, parallel I/O, performance analysis and tuning, power, programming models (data parallel, task parallel, process-centric, shared/distributed memory), scalability and performance studies, scheduling, storage systems, synchronization, and tools (Cuda, Swift, Globus, Condor, Amazon AWS, OpenStack, Cilk, gdb, threads, MPICH, OpenMP, Hadoop, FUSE). \par\vspace{2pt}{\footnotesize\itshape Set texts: Distributed Systems: Principles and Paradigms, A. S. Tanenbaum and M. V. Steen, Prentice Hall, 2nd Edition, 2007; Distributed and Cloud Computing: Clusters, Grids, Clouds, and the Future Internet, K Hwang, J Dongarra and GC. C. Fox, Elsevier, 1st Ed. 65 Course Contents of Mathematics \& Supporting Courses.}}
\edrow{2025}{3 \;$\bullet$\; no prerequisite and no outline published}{\begin{tight}
\item Explain concepts of concurrency, synchronization, and distributed architectures. (C2: Understand)
\item Apply parallel programming libraries to solve computational problems. (C3: Apply)
\item Analyze the performance trade-offs between sequential and parallel solutions. (C4: Analyze)
\end{tight}}
\end{coursecard}


\begin{notebox}[Reading these three columns as a drafting instruction]
For any course in this part the practical rule is the same. If HEC 2023 gives an
outline, use it --- it is the most recent detailed specification and it carries
prerequisites. If only HEC 2017 gives one, use that and modernise the tooling; the
structure is almost always still right, and it is the only detailed text that
exists for the Information Technology courses. If only HEC 2025 gives outcomes,
the department is writing the course from scratch and should say so in its own
course file, mapping each locally written topic to one of the published CLOs so
that accreditation has something to check against.
\end{notebox}
%=====================================================================
\section*{Sources}
\addcontentsline{toc}{section}{Sources}
%=====================================================================
\markboth{Sources}{}

Every statement in this document is drawn from the three booklets below. Page
references are to the PDF page, not the printed page number.

\begin{longtable}{@{}L{3.4cm}L{12.4cm}@{}}
\toprule
\rowcolor{primary!15}
\textbf{Edition} & \textbf{Document and the sections used} \\
\midrule
\endhead
\textbf{\color{ed2017}2017} &
\emph{Curriculum of BS \& MS (CS, SE \& IT)}, HEC, 171 pp.
BS-IT curriculum structure and study plan, pp.\ 27--33; BS-CS domain courses,
pp.\ 23--24; BS-SE elective list (the only mention of Cloud Computing), p.\ 38;
the 75-course bachelor list, pp.\ 69--70; course details for Artificial
Intelligence, Big Data Analytics, Computer Networks, Computer Vision, Natural
Language Processing, System and Network Administration, IT Infrastructure and
Virtual Systems and Services, pp.\ 71--130; MS Data Science programme,
pp.\ 47--50 \\
\rowcolor{lightbg}
\textbf{\color{ed2023}2023} &
\emph{BS Curriculum, Computing Disciplines}, HEC / NCEAC, 100 pp.
Generic computing structure, pp.\ 10--12; BSCS model and study plan, pp.\ 17--19;
BSAI, pp.\ 23--25; BSDS, pp.\ 26--28; BSIT model and study plan, pp.\ 41--43;
full list of BS courses, pp.\ 50--53; course contents, pp.\ 54--100 --- note the
scope statement at the head of section~7, which limits published contents to the
courses common to all computing degrees plus the AI, Data Science and Cyber
Security cluster tables (7.4--7.8) \\
\textbf{\color{ed2025}2025} &
\emph{HEC Curriculum of Computer Science}, HEC, 86 pp.
Programme structure and interdisciplinary courses, pp.\ 15--16; scheme of studies
for Semesters I--VIII, pp.\ 17--20; specializations and the guidelines for them,
pp.\ 20--21; Data Science cluster, pp.\ 21--22; Artificial Intelligence cluster,
p.\ 22; Information Technology cluster, p.\ 23; Network Infrastructure \& Cloud
Computing cluster, pp.\ 17--18 of the specialization listing; course learning
outcomes for the mandatory major courses, pp.\ 24--26, and for the specialization
courses, pp.\ 37--64; AI specialization MOOC and certification recommendations,
pp.\ 67--73 \\
\bottomrule
\end{longtable}

\vspace{6pt}
\begin{notebox}[On the two courses this document does not settle]
The 2023 booklet lists both \emph{Data Science} and \emph{Introduction to Data
Science} in its master course list, and gives a full specification only for the
second. Likewise it lists both \emph{Artificial Neural Networks \& Deep Learning}
and \emph{Deep Learning} with overlapping outlines. Where the two conflict this
document follows the specification with the published outline and prerequisites,
and says so in the relevant card.
\end{notebox}

\end{document}
